\documentclass[11pt]{article}
\usepackage[T1]{fontenc}
\usepackage{lmodern}
\usepackage{microtype}
\usepackage[margin=1in]{geometry}
\usepackage{amsmath,amssymb,amsthm,mathtools}
\usepackage{booktabs,longtable,array}
\usepackage{enumitem}
\usepackage{xcolor}
\usepackage{listings}
\usepackage{url}
\usepackage[hidelinks]{hyperref}
\hypersetup{pdftitle={CIFER: Certified Interfaces, Finite Evidence, and Residualized Proof Compilation for Intuitionistic PRA},pdfauthor={Research report prepared for Vladimir Reshetnikov},pdfsubject={A general framework for embedding finite mathematical proofs into iPRA}}
\usepackage[nameinlink,noabbrev]{cleveref}
\usepackage{fancyhdr}
\usepackage{lastpage}

\setlist{nosep,leftmargin=2em}
\setlength{\parindent}{0pt}
\setlength{\parskip}{0.55em}
\setlength{\emergencystretch}{2em}
\setlength{\headheight}{14pt}
\pagestyle{fancy}
\fancyhf{}
\lhead{CIFER for iPRA}
\rhead{\thepage/\pageref{LastPage}}
\cfoot{}

\definecolor{codegray}{gray}{0.96}
\lstset{
  basicstyle=\ttfamily\small,
  backgroundcolor=\color{codegray},
  frame=single,
  columns=fullflexible,
  breaklines=true,
  keepspaces=true,
  showstringspaces=false
}

\newtheorem{theorem}{Theorem}[section]
\newtheorem{lemma}[theorem]{Lemma}
\newtheorem{proposition}[theorem]{Proposition}
\newtheorem{corollary}[theorem]{Corollary}
\newtheorem{definition}[theorem]{Definition}
\newtheorem{remark}[theorem]{Remark}
\newtheorem{example}[theorem]{Example}
\newtheorem{warning}[theorem]{Warning}

\newcommand{\iPRA}{\ensuremath{\mathsf{iPRA}}}
\newcommand{\PRAAx}{\ensuremath{\mathsf{PRAAx}}}
\newcommand{\Prf}{\ensuremath{\mathsf{Prf}}}
\newcommand{\NPrf}{\ensuremath{\mathsf{NPrf}}}
\newcommand{\Nat}{\ensuremath{\mathbb{N}}}
\newcommand{\Bool}{\ensuremath{\mathbb{B}}}
\newcommand{\code}[1]{\texttt{#1}}
\newcommand{\CIFER}{\textsc{CIFER}}
\newcommand{\res}{\operatorname{Res}}
\newcommand{\Reach}{\operatorname{Reach}}
\newcommand{\height}{\operatorname{ht}}
\newcommand{\diam}{\operatorname{diam}}
\newcommand{\supp}{\operatorname{supp}}

\title{\textbf{CIFER: Certified Interfaces, Finite Evidence, and\\Residualized Proof Compilation for Intuitionistic PRA}}
\author{Research report prepared for Vladimir Reshetnikov}
\date{15 November 2026 (America/Los Angeles)}

\begin{document}
\maketitle

\begin{abstract}
The existing ProveIt development contains a kernel-checked presentation of intuitionistic primitive recursive arithmetic and a large collection of finitary proof packets.  The packets repeatedly discover the same successful ingredients -- fixed primitive-recursive codes, bounded search, quantitative names, finite witness clouds, backward precision propagation, and proof-carrying traces -- but most applications stop one level before a derivation of type \(\Prf\,\PRAAx\).  Moreover, several difficult theorems depend on genuinely missing quantitative producers: a height bound, covering representatives, an effective analytic constant, or a separation gap.

This report develops \CIFER, a framework that turns both facts into part of the formal statement.  A query is specialized to a finite interface graph.  Every executable node carries a local object-theory derivation constructor; every unavailable producer is represented by an explicit residual boundary formula.  A reverse dependency pass computes the canonical dependency-normal residual, a demand pass fixes all finite resources before induction, and a forward proof pass constructs
\[
   B_1,\ldots,B_r,\Gamma\ \vdash_{\PRAAx}\ G.
\]
Supplying embedded proofs of the \(B_i\) closes the theorem by simultaneous cut.  If no such proofs are available, the result remains a precise conditional theorem rather than an overstated solution.

The report gives full finitary proofs of residual compilation, fixed-formula Booleanization, least bounded search, capped cyclic demand closure, finite-trace reflection, backward viability, stable-observable elimination, shrinking finite-horizon coherence, and height/list conversion.  It then applies the framework to the recurring obstructions in Diophantine geometry, represented analysis, exact branch selection, and effective prime distribution.  A detailed Lean integration plan targets the existing declarations \code{PR}, \code{PTm}, \code{PRAAx}, \code{Prf.trans}, \code{PRA.qfInd}, and \code{pra\_iff\_native}.  The package includes uncompiled Lean source with no admitted theorems, exact Python reference algorithms, regression tests, and a worked rational bisection example for \(y^3+y=x\).
\end{abstract}

\tableofcontents
\newpage

\section{Executive result and use of the framework}

\subsection{The central problem}

A mathematical proof can fail to become a proof in \iPRA\ for at least four very different reasons.

First, the underlying theorem may be available in an ordinary mathematical language but its operations have not been compiled to primitive-recursive syntax.  Second, an algorithm may be primitive recursive and tested, while its invariant has not been turned into a quantifier-free object-language formula and proved by the induction rule of \iPRA.  Third, the mathematics may determine only a finite relation of acceptable approximate witnesses, while the attempted statement asks for one exact, continuously chosen witness.  Fourth, the proof may genuinely require quantitative information that has not been produced: a height cutoff, a modulus, a list of rational representatives, a zero-free region, or a separation gap.

These cases should not be conflated.  In particular, an executable host program is not an object-theory proof; a finite cardinality bound is not a search cutoff; a finite cloud is not a canonical branch; and an unproved producer should not be hidden inside a definition called ``data.''  The purpose of \CIFER\ is to enforce these distinctions while retaining as much of each proof as is already constructive and finitary.

\subsection{One output type instead of many informal statuses}

Fix a number \(n\) of free numerical variables and a permanent input context \(\Gamma\) of formulas over those variables.  A \CIFER\ compilation has the following formal endpoint:
\begin{equation}\label{eq:residual-endpoint}
 B_1,\ldots,B_r,\Gamma\ \vdash_{\PRAAx}\ G.
\end{equation}
Here \(G\) is the requested theorem and the formulas \(B_i\) are the residual producer obligations.  Every step from the \(B_i\) and \(\Gamma\) to \(G\) is intended to be represented by a concrete Lean value of type \code{Prf PRAAx}.  The residual list is computed from the dependency graph of the selected query, so unused global assumptions do not appear.

There are then three clean outcomes.

\begin{enumerate}
\item If \(r=0\), the result is a closed theorem of \iPRA.
\item If object-language derivations \(\Gamma\vdash B_i\) are later supplied, simultaneous cut closes \eqref{eq:residual-endpoint} without changing the downstream proof.
\item If some \(B_i\) remains unavailable, the package still contains a mathematically meaningful conditional theorem and an exact account of what is missing.
\end{enumerate}

This is more informative than a binary ``proved/not proved'' flag.  It also makes progress monotone: improving a producer removes a boundary node, while all verified consumers remain reusable.

\subsection{The framework in six passes}

A query-specific \CIFER\ build proceeds through six finite passes.

\paragraph{Pass 1: representation audit.}
Each ideal object is replaced by the finite information actually consumed by the theorem.  Exact integers and rationals use canonical codes.  Reals use precision-indexed rational enclosures or approximants.  Roots, eigenvectors, local branches, and geometric choices may use finite verified clouds instead of a global selector.  Quotient objects carry explicit equality certificates.

\paragraph{Pass 2: graph specialization.}
The host system constructs a fixed finite graph for the selected theorem.  It specializes all operations to concrete primitive-recursive syntax.  There is no object-theory interpreter for arbitrary encoded programs.  This avoids the impossible demand for a total primitive-recursive universal evaluator.

\paragraph{Pass 3: backward residual and demand analysis.}
Backward traversal identifies exactly the unresolved boundary leaves used by the goal.  A separate monotone closure computes precision, fuel, list, and search demands.  Cyclic precision dependencies are resolved on a finite capped lattice before any induction is invoked.

\paragraph{Pass 4: forward finite evidence.}
Primitive-recursive transducers, bounded loops, bounded searches, and finite relational viability tables are evaluated with the fixed demands.  This pass produces traces and certificates.  It is useful for testing, but it is not yet the proof.

\paragraph{Pass 5: object-proof composition.}
Each node carries a local derivation from its declared dependencies.  Quantifier-free loop invariants are proved by \code{PRA.qfInd}; local results are composed in topological order by \code{Prf.trans}.  The output is exactly \eqref{eq:residual-endpoint}.

\paragraph{Pass 6: closure or honest export.}
Available producer proofs discharge the residual context.  The resulting usual-iPRA derivation may then be transported through the checked \code{pra\_iff\_native} bridge.  Otherwise, the conditional result and residual report are exported together.

\subsection{What is new relative to the existing finite-interface proposals}

The existing repository contains several strong and mutually compatible ideas: certified finitary presentations, finite-use relations, proof-carrying observation graphs, query-specialized state machines, demand closure, stable observables, and local-to-global certificates.  \CIFER\ does not replace these ideas.  Its new organizing principle is that an unavailable producer is a first-class formula in the output sequent.

This residualization has five consequences.

\begin{enumerate}
\item \textbf{A generic compiler can be complete without pretending to solve deep mathematics.}  Completeness means that it derives the target from the exact unresolved boundaries, not that every boundary is automatically inhabited.
\item \textbf{Difficulty reports become compositional artifacts.}  A report can name a formula, verifier, and downstream theorem.  Later work needs only to prove that formula.
\item \textbf{Irrelevant assumptions disappear.}  Residual extraction is query-local.  A height bound for one family does not contaminate an unrelated finite algebra theorem merely because both live in the same library.
\item \textbf{Alternative constructive conclusions are explicit.}  If an exact selector needs a gap, the no-gap output can remain a finite relation, a stable observable, or an approximation theorem rather than being silently strengthened.
\item \textbf{Proof status is machine-auditable.}  A closed endpoint has type \code{Prf PRAAx [] goal}; a conditional endpoint has a nonempty boundary context; a host test has neither type.
\end{enumerate}

\subsection{Main generic theorems}

The human arguments in this report establish the following generic results.

\begin{description}[style=nextline]
\item[Residual compilation.] Every finite acyclic proof graph with local embedded derivations yields a derivation of the target from the dependency-normal set of reachable boundary leaves.
\item[Fixed bounded reflection.] Every fixed quantifier-free or strictly bounded formula can be compiled to a primitive-recursive Boolean checker, with truth and falsity reflected in \iPRA.
\item[Least bounded search.] A bounded scan returns the least witness and proves soundness, minimality, and the no-witness case by one quantifier-free invariant.
\item[Capped demand closure.] An inflationary monotone precision transformer on a finite product of capped natural orders reaches its least fixed point after an explicit number of strict rank increases.
\item[Finite-trace reflection.] A fixed primitive-recursive trace with a quantifier-free invariant yields the invariant at arbitrary primitive-recursive fuel and an unrestricted final formula through a terminal extractor.
\item[Finite relational viability.] Backward dynamic programming decides whether a finite compatible path exists and constructs the lexicographically least path without invoking infinite choice.
\item[Stable-observable elimination.] If all valid witnesses have the same finite observation, or observations lie in an explicitly small diameter, the theorem can eliminate internal witness choice.
\item[Shrinking finite-horizon coherence.] Least viable choices at changing horizons form a Cauchy name under a shrinking-diameter modulus and a summable motion budget; no infinite branch object is required.
\item[Height/list conversion.] For a decidable coded predicate, a height bound yields a complete finite list, and a complete finite list yields a height bound.  A mere cardinality bound supplies neither a stopping criterion nor the missing codes.
\end{description}

\subsection{Scope of the claimed resolution}

The report resolves architectural and logical difficulties, not every theorem-specific mathematical difficulty.  For example, once a correct height bound is supplied, the conversion to complete finite enumeration is routine, uniform, and formalizable.  The framework does not manufacture a general effective Faltings or Siegel bound.  Once rational covering representatives for a Mordell--Weil descent are supplied, the finite descent and generator extraction are routine.  The framework does not turn Selmer classes into rational representatives without proof.  Once explicit analytic constants are supplied, finite prime-sum estimates can be compiled.  The framework does not conjure a zero-free region or an effective exceptional-zero bound.

This boundary is a feature.  A general framework should solve the recurring proof-engineering problem completely and expose irreducible theorem-specific inputs exactly.

\subsection{Artifact status}

The generic arguments are fully written at the mathematical level.  The accompanying Python implementation is an exact finite reference and its regression tests are executed in this package.  The Lean files use the present repository names and include an object-language composition core, host graph definitions, demand and relational blueprints, and application interfaces.  They contain no admitted theorem, but the package does not claim that they were elaborated: no Lean executable was available in the artifact runtime.  Consequently, no new theorem in this report is labeled kernel-checked.

\section{Repository audit and recurring failure modes}

\subsection{The checked foundation and the unverified packet layer}

The audited repository has two sharply different layers.

The first is \code{Logic/PrimitiveRecursiveArithmetic/ConciseIPRA}.  It contains a deeply embedded syntax, intuitionistic natural deduction with equality, primitive-recursive function symbols, the usual \iPRA\ axioms, a native scalar calculus, and checked translations between the two presentations.  The declarations relevant to this report are:
\begin{itemize}
\item \code{Fm}, substitution, and \code{Fm.QF} in \code{Syntax.lean};
\item \code{Prf}, \code{Prf.trans}, substitution, weakening, and equality transport in \code{Deduction.lean};
\item \code{PR}, \code{PTm}, \code{PRAAx}, and \code{PRA.qfInd} in \code{PR.lean};
\item \code{pra\_iff\_native} and \code{native\_iff\_pra} in \code{NativeBiInterpretation.lean}.
\end{itemize}
These declarations give a precise endpoint.  A completed theorem is not merely a Lean proposition about natural numbers; it is an inhabitant of an embedded derivability judgment.

The second layer is \code{docs/proofs}.  At the audited snapshot it catalogs seventy-eight delivered packets across algebra, analysis, number theory, geometry, probability, topology, and foundations.  The repository README explicitly records that these packets were placed as delivered and were generally not rerun, compiled, or edited.  Many Lean files are honest blueprints; some native finite cores compile; a small number of framework combinators construct actual \code{Prf PRAAx} values.  The main mathematical packets, however, do not yet end in checked embedded derivations of their headline theorems.

That distinction determines the proper research task.  It is not enough to invent another semantic interface at the host level.  The missing bridge is a repeatable route from a query-specific finite construction to an actual object-language derivation, while retaining explicit boundaries for mathematical producers not yet available.

\subsection{Convergence among the existing framework proposals}

The foundation packets use different terminology, but their strongest ideas converge on a common architecture.

\begin{longtable}{p{0.22\textwidth}p{0.69\textwidth}}
\toprule
Recurring idea & Typical role \\
\midrule
\endhead
Fixed finite presentation & Replace an ideal object by codes and finite observations used by one query. \\
Primitive-recursive checker & Turn a certificate into a decidable, bounded arithmetic relation. \\
Quantifier-free invariant & Justify a bounded trace using precisely the induction available in \iPRA. \\
Backward demand & Determine required precision, fuel, or list size before running a producer. \\
Relational witness cloud & Preserve nonunique or collision-prone approximate witnesses without forced selection. \\
Stable observable & Eliminate witness choice when all valid branches give the same requested output. \\
Finite viability & Remove dead-end approximate witnesses at a fixed horizon by backward dynamic programming. \\
Proof graph & Compose local derivations in dependency order. \\
Producer/verifier split & Isolate deep existence or bound information from elementary finite checking. \\
Query specialization & Compile concrete \code{PR} syntax instead of postulating a universal evaluator. \\
\bottomrule
\end{longtable}

The Proof-Carrying Observation Graph packet comes closest to a comprehensive synthesis.  OCiPRA gives a useful classification into deterministic contracts, bounded relational contracts, stable observables, and conditional theorems with explicit producer data.  The kernel finite-use packet gives source expressions for topological proof composition, bounded loops, stable observations, explicit boundaries, and export through \code{pra\_iff\_native}.  The PFIG packet emphasizes normalization of a fixed graph to one query-specific state machine.  Other packets contribute coherent finite observations, finite demand, local-to-global compilation, categorical gluing, and exact application adapters.

The main opportunity is therefore consolidation around the checked endpoint, not proliferation of names.

\subsection{Failure mode A: host truth in place of object derivability}

A Lean theorem of type
\[
   \forall x:\Nat,\ P(x)
\]
can be valuable evidence about an algorithm, yet it is not automatically a theorem of the embedded arithmetic.  To obtain the latter, one must provide a formula \(\widehat P\), primitive-recursive terms for every operation, and a term of type
\[
  []\ \vdash_{\PRAAx}\ \forall x\,\widehat P(x).
\]
The gap cannot be closed by a cast.  The syntax and proof theory are data inside Lean, and their derivations must be constructed explicitly.

This issue recurs whenever a packet proves a native theorem about lists, rational arithmetic, or a finite search and then supplies only a conditional bridge saying that an already constructed embedded derivation can be translated.  \CIFER\ resolves the architectural part by requiring every proof-graph node to carry an object-theory local derivation constructor.  Host computation may generate syntax and evidence, but it does not certify itself.

\subsection{Failure mode B: hidden precision or fuel choices}

An informal effective proof often says ``choose the input precision sufficiently large'' or ``iterate until the error is small.''  Such phrases hide an existential choice that may depend cyclically on other choices.  In \iPRA, all loop bounds must be represented by primitive-recursive terms, and the induction formula must be fixed before induction starts.

The existing demand-passing proposals correctly move this choice into a reverse pass.  The remaining complication is feedback: an output tolerance may demand an input precision, which changes a denominator bound, which in turn changes the output tolerance.  \CIFER\ treats the demand vector as a point in a finite capped product order and computes an inflationary fixed point.  The explicit cap is part of the theorem interface; stabilization is proved by a rank argument.

\subsection{Failure mode C: cardinality confused with search termination}

Several Diophantine packets correctly emphasize that knowing there are at most \(N\) solutions does not tell a bounded search when all solutions have been found.  The distinction is simple but fundamental.

A height bound \(H\) supplies a finite search domain.  A complete finite list supplies a height bound by taking a maximum.  A cardinality bound alone supplies neither a domain nor a completeness certificate.  Searching larger and larger boxes until \(N\) points have appeared is unsound when the true number is smaller than \(N\), and even when the true number equals \(N\), the equality itself may not be known constructively.

\CIFER\ therefore uses separate formula types for \emph{count}, \emph{cutoff}, and \emph{complete list}.  Only the latter two are interchangeable under decidability.

\subsection{Failure mode D: exact selectors through collisions}

Many classical statements hide a discontinuous choice.  Examples include an axis for the identity rotation, one ordered root through a multiple-root collision, a normalized eigenvector through an eigenvalue collision, or one local analytic branch at a singular point.  A Cauchy name for the input need not determine a stable exact label.

A common but invalid repair is to choose the lexicographically least approximate witness at each precision and call the resulting sequence the exact witness.  Tie-breaking can jump forever.  The correct constructive outputs are one of the following:
\begin{enumerate}
\item a finite relation of valid witnesses at each requested precision;
\item an exact observable invariant across all those witnesses;
\item an approximate observable with an explicit diameter bound;
\item an exact label under a supplied positive separation gap.
\end{enumerate}
The framework makes the chosen output class part of the statement, not an implementation detail.

\subsection{Failure mode E: infinite coherence demanded when only finite use is needed}

A theorem consuming a represented real at precision \(q\) reads only finitely much of its name.  Nevertheless, formalizations often begin by constructing a globally coherent infinite witness sequence.  This introduces unnecessary choice and coherence obligations.

The finite-use principle reverses the order.  Fix the query and its demanded precision.  Construct a finite compatible path to the corresponding horizon.  If the consumer observes only the endpoint or a stable quotient, no infinite object is needed.  A coherent Cauchy name is constructed only when the theorem itself requests one; even then, changing finite horizons plus explicit shrinking estimates can replace an actual infinite branch choice.

\subsection{Failure mode F: a universal evaluator smuggled into PRA}

The class of primitive-recursive functions is effectively enumerable, but there is no total primitive-recursive function that evaluates every primitive-recursive program on every input.  A universal evaluator for all programs would diagonalize out of the class.  Consequently, a framework must not put an unrestricted interpreter inside \iPRA\ and then appeal to its correctness.

The repository's better proposals already avoid this.  A host compiler sees a fixed finite program or proof graph and emits its concrete \code{PR} syntax.  All subsequent induction concerns that fixed syntax.  \CIFER\ adopts this query-local specialization as a nonnegotiable design rule.

\subsection{Failure mode G: genuinely missing quantitative producers}

Some gaps are not proof-engineering gaps.  They are missing mathematical data.

\begin{itemize}
\item A general effective Siegel or Faltings application needs a positive height bound, not merely a finiteness or cardinality statement.
\item A constructive Mordell--Weil descent needs actual rational representatives satisfying a covering assertion; finite Selmer information alone does not automatically provide them.
\item Effective prime-distribution theorems need explicit field or conductor parameters, zero-free regions, density constants, or exceptional-zero control.
\item Exact branch recovery needs a separation certificate when the input presentation permits collisions.
\end{itemize}

A general compiler cannot create these producers.  It can, however, prove every downstream consequence uniformly and return the producer formulas as residual obligations.  This is the key role of residualized compilation.

\subsection{Failure mode H: quotient equality and local-to-global gluing}

Finite presentations of algebraic and geometric objects frequently use noncanonical representatives.  Equality in a localization, ideal quotient, projective coordinate system, simplicial homology class, or sheaf section is witnessed by finite algebraic data.  Treating representation equality as literal code equality destroys extensionality; quotienting at the host type level may hide the witness needed by \iPRA.

The robust pattern is a \emph{certified setoid}: codes, a primitive-recursive equality verifier, reflexivity/symmetry/transitivity constructors, and extensional operation proofs.  Gluing is then a finite compatibility theorem.  \CIFER\ uses certified setoids as one of its representation adapters and treats the equality certificate as ordinary finite evidence.

\subsection{Audit conclusion}

The repository already contains most ingredients of a successful methodology.  The remaining general task is to make the endpoint and boundaries unavoidable.  Every application should answer five questions:

\begin{enumerate}
\item What finite observation does the query actually consume?
\item Which algorithms and verifiers are fixed primitive-recursive syntax?
\item Which loops have explicit quantifier-free invariants and fuel?
\item Which local conclusions are actual \code{Prf PRAAx} derivations?
\item Which unavailable producers remain as formulas in the residual context?
\end{enumerate}

The rest of the report develops a calculus in which these questions are represented directly.

\section{Logical target: what must be internalized}

\subsection{The ordinary presentation of iPRA}

For the purposes of this report, \iPRA\ is first-order intuitionistic arithmetic with one numerical domain, one function symbol for every primitive-recursive function, the defining equations of those symbols, the usual successor axioms, and induction restricted to quantifier-free formulas.  In the repository, primitive-recursive symbols of arity \(k\) are represented by \code{PR k}; terms over \(n\) free variables by \code{PTm n}; and formulas by \code{Fm PTm n}.  The theory \code{PRAAx} contains projection, composition, and primitive-recursion equations, successor nonzeroness and injectivity, and the quantifier-free induction scheme.

A sequent
\[
 \Gamma\vdash_{\PRAAx} A
\]
is the Lean proposition \code{Prf PRAAx Gamma A}.  The proof object is built from intuitionistic natural deduction, equality, theory axioms, and substitution.  The central structural rule for compilation is simultaneous cut:
\[
 \frac{\Gamma\vdash A \qquad \text{for every }B\in\Gamma,\ \Delta\vdash B}
      {\Delta\vdash A}.
\]
This is exposed as \code{Prf.trans}.  It is exactly what a topological proof compiler needs: a node's local hypotheses can be replaced by the already compiled derivations of its dependencies.

\subsection{The native scalar presentation}

The ConciseIPRA development also contains a much smaller native syntax based on scalar iteration and formulas generated by equations and two binders.  The checked theorems \code{native\_iff\_pra} and \code{pra\_iff\_native} show that the native calculus and the ordinary presentation derive corresponding sequents.  This has two consequences for the framework.

First, it is enough to build object proofs in the ordinary presentation, whose connectives and quantifiers are convenient for application statements.  A finished proof can be exported to the native calculus.  Second, the bridge is not a substitute for the proof: \code{pra\_iff\_native} translates an existing derivation; it does not create one from a semantic theorem.

\subsection{Fixed formulas rather than internal syntax interpretation}

The metatheorems below repeatedly use the phrase ``for every fixed formula.''  This qualification is essential.

Let \(A(\vec x)\) be a fixed quantifier-free formula assembled from equations of primitive-recursive terms by propositional connectives.  At the host level, recurse over the finite syntax tree of \(A\).  For each node, construct a concrete primitive-recursive Boolean term and concrete local derivations.  The resulting Booleanizer is a primitive-recursive function because the formula is fixed.

This is not the same as defining one primitive-recursive function
\[
  \operatorname{Eval}(\ulcorner A\urcorner,\vec x)
\]
that evaluates arbitrary coded primitive-recursive terms.  Such a universal total evaluator cannot itself be primitive recursive.  \CIFER\ always performs syntax recursion before entering the object theory.

\begin{remark}
The host compiler may be general software.  Generality at the host level is harmless because each invocation emits a finite, specialized object program and proof.  The object theory sees only that specialization.
\end{remark}

\subsection{The finite evidence discipline}

Every mathematical notion used by a compiled theorem is divided into four layers.

\begin{definition}[Code]
A code is a natural number or a fixed tuple of natural numbers represented by a primitive-recursive pairing scheme.  All operations on codes used by a theorem have concrete \code{PR} symbols.
\end{definition}

\begin{definition}[Verifier]
A verifier is a fixed primitive-recursive Boolean predicate on codes.  Its object formula is usually an equation \(v(\vec x)=1\), hence quantifier-free.
\end{definition}

\begin{definition}[Evidence]
Evidence is a finite code whose verifier succeeds: a trace, factorization, interval, finite list, compatibility table, equality certificate, or local proof certificate.
\end{definition}

\begin{definition}[Producer]
A producer is a primitive-recursive operation that returns evidence from declared inputs and demands.  When no such operation is available, the desired evidence is represented by a boundary formula rather than by an opaque host constant.
\end{definition}

This discipline separates two proof obligations.  \emph{Verification} proves that a supplied code has the advertised mathematical meaning.  \emph{Production} proves that an appropriate code exists and is obtained by a primitive-recursive operation.  Verification is often elementary even when production is deep.

\subsection{Quantifier-free induction as a trace rule}

Suppose \(S(t,\vec x)\) is the coded state after \(t\) iterations of a fixed primitive-recursive transition, and \(I(t,S(t,\vec x),\vec x)\) is quantifier-free.  The ordinary proof pattern is:
\begin{align*}
 &I(0,S(0,\vec x),\vec x),\\
 &I(t,S(t,\vec x),\vec x)\to I(t+1,S(t+1,\vec x),\vec x),\\
 &I(T,S(T,\vec x),\vec x)\to G(\vec x).
\end{align*}
The first two lines are exactly the premises of \code{PRA.qfInd}; universal elimination specializes the induction theorem at the primitive-recursive fuel \(T\); and the last line may conclude an arbitrary formula \(G\).  Thus the restriction to quantifier-free induction does not require the final theorem to be quantifier-free.  It requires the iterative state invariant to be so.

This observation supports a general design rule:

\begin{quote}
Put all unbounded logical structure outside the loop.  Inside the loop, carry finite Boolean facts, equalities, inequalities, bounded counters, and certificate accumulators.
\end{quote}

\subsection{Bounded quantifiers are finite programs}

A strictly bounded existential
\[
  \exists z<b\;R(z,\vec x)
\]
with decidable \(R\) is represented by a finite disjunction or a Boolean fold.  A bounded universal is represented by a finite conjunction.  The fold is primitive recursive in \(b\), and its reflection theorem is proved by quantifier-free induction on the fold length.  Bounded minimization similarly scans \(0,1,\ldots,b-1\), maintaining a flag and the least witness found.

Accordingly, bounded quantifiers may appear in the human specification without increasing the logical strength of the iterative core.  The compiler expands them into fixed primitive-recursive checkers.

\subsection{Coding tuples, lists, and finite maps}

All finite structures can be reduced to one numerical sort.  For a fixed application, choose one pairing function \(\langle a,b\rangle\) with primitive-recursive projections.  A tuple of fixed dimension is nested pairing.  A list is encoded by length plus entries, or by a bounded sequence code.  A finite Boolean table of size \(m\times n\) is encoded by a bit vector or by a sequence of entries.

The framework does not depend on a particular code, but a formal application must provide the following equations and bounds:
\begin{enumerate}
\item projection after construction;
\item update at the selected index and preservation elsewhere;
\item bounded access default behavior;
\item length and range bounds;
\item extensionality under pointwise equality on the declared range.
\end{enumerate}
These are elementary quantifier-free theorems and should live in one reusable arithmetic interface rather than be reproved by every domain packet.

\subsection{Certified setoids for quotient-like objects}

Literal equality of codes is often too strong.  A certified setoid presentation consists of:
\[
 (C,\;\operatorname{Valid},\;\operatorname{EqCert},\;\operatorname{checkEq}),
\]
where \(C\) is the code type, \(\operatorname{Valid}(x)\) is a Boolean validity condition, and \(\operatorname{checkEq}(x,y,e)=1\) verifies a finite certificate that \(x\) and \(y\) represent the same object.  The application supplies primitive-recursive constructors proving reflexivity, symmetry, and transitivity, together with congruence of operations.

Examples include:
\begin{itemize}
\item rationals as signed numerator/positive denominator pairs with cross-multiplication certificates;
\item projective points with a nonzero scaling relation;
\item localization equality witnessed by a killing exponent;
\item ideal membership witnessed by a finite linear combination;
\item homology equality witnessed by a boundary chain;
\item finite atlases whose overlap equality is witnessed chart by chart.
\end{itemize}

This avoids both extremes: one need not quotient the host type and lose explicit witnesses, and one need not mistake representation equality for mathematical equality.

\subsection{Statement classes}

A mature application should identify its theorem class before proving it.

\begin{description}[style=nextline]
\item[Deterministic theorem.] A primitive-recursive producer returns one code and the verifier proves the exact result.
\item[Bounded relational theorem.] A producer returns a nonempty finite set or relation of valid outputs.
\item[Stable-observable theorem.] The internal relation is nonunique, but every valid output gives the same requested finite observation.
\item[Approximate theorem.] All valid outputs lie in a declared diameter at the requested precision.
\item[Gap-conditional exact theorem.] Approximate evidence plus a positive separation certificate determines one exact label.
\item[Residual conditional theorem.] Downstream derivation is complete, but one or more producer formulas remain in the context.
\end{description}

These classes are not ranked by prestige.  The correct class is the strongest one supported by the available data without importing a hidden choice principle.

\subsection{The formal contract of a node}

A proof-graph node \(v\) has:
\begin{enumerate}
\item a formula \(A_v\);
\item an ordered list of predecessor nodes \(u_1,\ldots,u_k\);
\item a kind: input, local, boundary, or final;
\item for a local or final node, a derivation
\[
 A_{u_1},\ldots,A_{u_k},\Gamma\vdash_{\PRAAx} A_v;
\]
\item optional executable metadata: a producer, verifier, demand transformer, and trace schema.
\end{enumerate}
An input formula is already in \(\Gamma\).  A boundary formula has no local derivation.  The graph is topologically ordered at the proof level.  Precision-demand metadata may contain cycles, but those cycles are solved separately before forward proof composition.

This node contract is the smallest unit at which host generation, finite execution, and object proof can be audited independently.

\section{CIFER interfaces and query-local compilation}

\subsection{Finite observation interfaces}

The same ideal mathematical object may support many finite interfaces.  A continuous real function, for example, can be consumed through pointwise values, interval enclosures, a modulus on a compact interval, a finite graph approximation, or polynomial coefficients.  A universal representation is rarely the simplest route to \iPRA.  \CIFER\ therefore makes the interface relative to a query.

\begin{definition}[Finite observation interface]
A finite observation interface for inputs \(x\) and demand \(q\) consists of:
\begin{enumerate}
\item an evidence code \(e\);
\item a quantifier-free verifier \(V(x,q,e)\);
\item an observation code \(o=\operatorname{obs}(e)\);
\item a correctness formula \(C(x,q,o)\);
\item an embedded derivation
\[
 V(x,q,e)\vdash_{\PRAAx} C(x,q,\operatorname{obs}(e)).
\]
\end{enumerate}
The interface may additionally contain a primitive-recursive producer \(E(x,q)\) and a proof of \(V(x,q,E(x,q))\).  If the producer is unavailable, the existence or validity formula becomes a residual boundary.
\end{definition}

The observation may itself be relational.  In that case \(e\) codes a finite list \(L\), the verifier proves that every member is valid and that the list is nonempty or complete for the finite query, and the consumer works with all members or with a stable observable.

\subsection{Demand transformers}

Every interface declares what input information is needed for an output demand.  For a deterministic operation \(f\), a demand transformer is a primitive-recursive function
\[
 d_f(q,\alpha)=p,
\]
where \(q\) is the requested output precision and \(\alpha\) collects finite side information such as denominator exponents, coefficient heights, or a supplied gap.  The local correctness theorem has the form
\[
  C_X(x,p,a)\to C_Y(f(x),q,F(a)).
\]

For a graph node with several inputs, the transformer returns one demand per predecessor.  The reverse pass takes componentwise maxima over all consumers.  When a transformer's side information depends on demands elsewhere, the graph gives a monotone system rather than a simple acyclic recurrence.  Section~\ref{sec:demand-closure} solves that system under an explicit cap.

\subsection{Evidence classes}

The framework uses a small number of evidence patterns repeatedly.

\paragraph{Exact calculation evidence.}
An exact node computes a primitive-recursive term and proves its defining equation.  Examples are integer polynomial evaluation, rational arithmetic, finite sums, matrix multiplication in fixed dimension, and list filtering.

\paragraph{Trace evidence.}
A trace node codes an initial state, a transition, a fuel, and an invariant.  Its public result is extracted from the terminal state.  Bounded search, Euclidean algorithms, descent, dynamic programming, finite Gaussian elimination, and interval iteration all fit this pattern.

\paragraph{Certificate evidence.}
A certificate need not be generated by the verifier.  Examples include a factorization, a Bezout identity, an ideal-combination witness, a nonzero resultant, a graph coloring, a finite cover, or a rational point.  The verifier is primitive recursive and the theorem says what follows if it accepts.

\paragraph{Relational evidence.}
A relational node returns a finite cloud of candidates together with validity and compatibility tables.  It supports collision-prone or noncanonical data.

\paragraph{Observation evidence.}
An observation node proves that all valid candidates have a common label, lie in a small interval, induce the same homology map, or yield the same bounded consumer result.

\paragraph{Boundary evidence.}
A boundary node states a producer result that is used downstream but not constructed in the current package.  Examples are ``this code is a complete height bound,'' ``these are rational covering representatives,'' and ``this rational is an effective zero-free constant.''

\subsection{Proof dependencies versus data dependencies}

It is useful to distinguish three graphs.

\begin{enumerate}
\item The \emph{proof graph} is a finite DAG.  A local proof may use only earlier nodes and the permanent context.
\item The \emph{data-flow graph} records which computed codes feed which operations.  It is usually a DAG after bounded loops are packaged as trace nodes.
\item The \emph{demand graph} records how requested resources propagate.  It may contain cycles, because a bound can depend on a precision that depends on the same bound.
\end{enumerate}

Conflating these graphs leads to confusion.  A demand cycle does not make the proof circular if it is resolved by a finite fixed-point computation before the proof run.  Conversely, a topologically ordered executable computation does not justify a proof node unless the local object derivation is supplied.

\subsection{Query specialization}

A query consists of:
\begin{itemize}
\item a target formula \(G\);
\item permanent input formulas \(\Gamma\);
\item a requested finite observation or precision;
\item an application library of possible interfaces.
\end{itemize}
The host compiler selects only the interfaces reachable from \(G\), specializes their terms and formulas, computes residual boundaries and demands, and emits a concrete proof schedule.

This specialization serves three purposes.

First, it limits the amount of formal arithmetic.  A theorem about a polynomial of fixed degree needs only fixed-dimensional tuple operations, not an internal general-purpose algebra system.  Second, it makes resource bounds sharper: the demand analysis sees the actual consumer.  Third, it avoids universal evaluation.  The generated object program is a fixed composition of \code{PR} constructors.

\subsection{A schema for deterministic contracts}

Let \(P(x)\) be a precondition, \(f(x)\) a fixed primitive-recursive output term, and \(Q(x,y)\) a postcondition.  A deterministic contract consists of
\[
 P(x)\vdash Q(x,f(x)).
\]
Composition is direct cut.  If
\[
 P\vdash Q(f),\qquad Q(y)\vdash R(g(y)),
\]
then instantiate the second local derivation at \(f\) and cut the first proof into it.  On a finite DAG, repeated simultaneous cut composes all nodes.

A deterministic node is appropriate only when the output really is a total primitive-recursive function of the declared input.  It should not be used to hide an unbounded search or a discontinuous selector.

\subsection{A schema for relational contracts}

A relational contract has a finite output list \(L(x,q)\) and a verifier proving
\[
 \operatorname{Nonempty}(L)\wedge
 \forall i<|L|\;R(x,q,L_i).
\]
If completeness is needed, add
\[
 \forall y\;\bigl(R^*(x,q,y)\to \exists i<|L|\;E(y,L_i,q)\bigr),
\]
where the quantification over ideal objects is replaced in the actual application by its finite coded or observational form.  The consumer may:
\begin{enumerate}
\item search the finite list for an additional decidable property;
\item propagate all candidates to another finite relation;
\item prune candidates by backward viability;
\item apply a stable observer;
\item require a gap and collapse to an exact label.
\end{enumerate}

The relation is therefore compositional.  Nonuniqueness is not a defect to be repaired prematurely.

\subsection{A schema for conditional producers}

Suppose \(B(x,b)\) says that \(b\) is a valid producer output and the downstream finite theorem is
\[
 B(x,b)\vdash G(x,b).
\]
There are three legitimate package levels.

\begin{enumerate}
\item \textbf{Verifier level:} prove the implication for arbitrary supplied \(b\).
\item \textbf{Producer level:} define a primitive-recursive \(p(x)\) and prove \(B(x,p(x))\).
\item \textbf{Closed level:} cut the producer proof into the verifier proof.
\end{enumerate}

When only the first is available, \CIFER\ reports \(B\) as a residual.  It does not introduce a noncomputable host constant \(b\), mark it opaque, and then present the result as a theorem of \iPRA.

\subsection{Residual normalization levels}

There are several notions of ``minimal boundary,'' and they should not be confused.

\begin{description}[style=nextline]
\item[Reachability normalization.] Remove boundary nodes not used by the target.
\item[Duplicate normalization.] Merge repeated occurrences of the same boundary node.
\item[Dependency expansion.] Replace any derived pseudo-boundary by the primitive unresolved leaves needed to derive it.
\item[Declared dominance.] If the library contains an embedded proof that one boundary follows from others, use it as an ordinary local node and eliminate the dominated boundary.
\item[Logical minimization.] Find a smallest logically sufficient subset of arbitrary formulas.  This is not a computable generic service and is not claimed.
\end{description}

The canonical \CIFER\ residual is minimal only in the first three syntactic senses, plus any explicitly supplied dominance proofs.  This level is deterministic, auditable, and sufficient for modular work.

\subsection{Proof-carrying interface records}

At the Lean level, an interface should not be a structure containing only semantic fields.  It should contain, or generate, a local embedded derivation.  Schematically:

\begin{lstlisting}
structure LocalNode (Gamma : List (Fm PTm n)) where
  formula : Fm PTm n
  deps : List (Fm PTm n)
  proof : deps ++ Gamma |- [PRAAx] formula
\end{lstlisting}

Executable metadata may be stored alongside the proof, but the proof field is decisive.  An application library can then expose two APIs:
\begin{enumerate}
\item a high-level builder that computes syntax and evidence;
\item a small trusted-facing record whose final field has type \code{Prf PRAAx}.
\end{enumerate}

This is the proof-theoretic analogue of proof-carrying code: the consumer checks the derivation rather than trusting the generator.

\subsection{Migration rule for existing packets}

An existing proof packet can be migrated incrementally.

\begin{enumerate}
\item Identify its headline target formula.
\item List all inputs currently described by phrases such as ``choose,'' ``there exists an effective,'' ``for a sufficiently large,'' or ``given a complete list.''
\item Classify each as a finite producer, a boundary, a relation, or a gap.
\item Turn every executable loop into a trace node with one explicit invariant.
\item Turn native Lean helper theorems into specifications for local object derivations, not into assumed bridges.
\item Compile the downstream DAG and record the residual context.
\item Close one small benchmark end to end before extending the domain library.
\end{enumerate}

This strategy preserves existing mathematical work while making proof status progressively more precise.

\section{Residual compilation theorem}
\label{sec:residual-compilation}

\subsection{Finite proof graphs}

Fix a permanent context \(\Gamma\).  A finite proof graph is a sequence of nodes
\[
 v_0,v_1,\ldots,v_{m-1}
\]
in topological order.  Node \(v_i\) carries a formula \(A_i\) and has one of three relevant kinds:

\begin{enumerate}
\item \emph{input}: \(A_i\in\Gamma\);
\item \emph{boundary}: no local derivation is supplied;
\item \emph{derived}: a finite predecessor list \(D_i\subseteq\{0,\ldots,i-1\}\) and a local derivation
\begin{equation}\label{eq:local-node-proof}
 (A_j)_{j\in D_i},\Gamma\vdash_{\PRAAx} A_i.
\end{equation}
\end{enumerate}
A final node is simply a designated derived or boundary node whose formula is the query target.  A boundary final is permitted: it states that the query itself remains an unresolved producer.

For clarity, predecessor lists are ordered and may contain repeated indices in raw input.  Normalization removes repetitions.  The proof does not depend on the order except for the exact list representation of the context.

\subsection{Recursive residual sets}

Define the dependency residual \(R(i)\) recursively along the topological order:
\[
 R(i)=
 \begin{cases}
  \varnothing,&v_i\text{ is input},\\
  \{i\},&v_i\text{ is boundary},\\
  \displaystyle\bigcup_{j\in D_i}R(j),&v_i\text{ is derived}.
 \end{cases}
\]
The union is stable and duplicate-free in graph order.  Let \(\mathcal B_i\) be the list of formulas \(A_j\) with \(j\in R(i)\).

\begin{remark}
The set \(R(i)\) is a syntactic dependency normal form.  It removes unreachable and internally derived assumptions.  It is not claimed to decide whether one arbitrary boundary formula is intuitionistically derivable from the others.
\end{remark}

\subsection{Main theorem}

\begin{theorem}[Residual compilation]\label{thm:residual-compilation}
For every node \(v_i\) of a finite proof graph,
\begin{equation}\label{eq:residual-node-sequent}
 \mathcal B_i,\Gamma\vdash_{\PRAAx} A_i.
\end{equation}
Consequently, if \(t\) is the target node, the graph compiles to
\[
 \mathcal B_t,\Gamma\vdash_{\PRAAx} A_t.
\]
\end{theorem}

\begin{proof}
Proceed by induction on \(i\).

If \(v_i\) is an input node, then \(R(i)=\varnothing\), and \(A_i\in\Gamma\).  Thus \eqref{eq:residual-node-sequent} is the hypothesis rule.

If \(v_i\) is a boundary node, then \(\mathcal B_i=[A_i]\).  Again the conclusion follows by the hypothesis rule.

Suppose \(v_i\) is derived.  For each \(j\in D_i\), the induction hypothesis gives
\[
 \mathcal B_j,\Gamma\vdash A_j.
\]
By weakening, each of these derivations is valid in the common context \(\mathcal B_i,\Gamma\), because \(R(j)\subseteq R(i)\).  The local proof \eqref{eq:local-node-proof} is a derivation from the predecessor formulas and \(\Gamma\).  Apply simultaneous cut, replacing every predecessor hypothesis \(A_j\) by its induction-hypothesis derivation and every permanent hypothesis by the hypothesis rule.  This yields
\[
 \mathcal B_i,\Gamma\vdash A_i.
\]
The three cases exhaust the node kinds.
\end{proof}

The proof is entirely intuitionistic and uses no semantic completeness theorem.  Its Lean realization is a fold of \code{ResidualBuild.stepFrom}; simultaneous cut is \code{Prf.trans}.

\subsection{Closure theorem}

\begin{corollary}[Boundary discharge]\label{cor:boundary-discharge}
Let the target residual be \(\mathcal B_t=[B_1,\ldots,B_r]\).  If
\[
 \Gamma\vdash B_k
\]
for every \(k\le r\), then \(\Gamma\vdash A_t\).
\end{corollary}

\begin{proof}
Apply simultaneous cut to the target derivation of \cref{thm:residual-compilation}, replacing each boundary hypothesis by its supplied derivation and each member of \(\Gamma\) by the hypothesis rule.
\end{proof}

This corollary is the modular closure point.  A downstream packet need not be regenerated when a producer is improved.  It is enough to supply a proof of the boundary formula with the same interface.

\subsection{Curried form}

Natural-deduction introduction yields a closed implication form:
\[
 \Gamma\vdash B_1\to(B_2\to\cdots(B_r\to A_t)\cdots).
\]
The context form is preferable for proof-graph composition because dependencies are unordered resources, while the curried form is convenient for statement export.  Both are intuitionistically equivalent by repeated implication introduction and elimination.

\subsection{Correctness of backward extraction}

The executable extractor starts at a target \(t\), marks it, and repeatedly marks every predecessor of a marked node.  After at most \(m\) sweeps, no new node can be marked.  On a topologically ordered graph, a single descending sweep suffices, but the repeated form is easier to implement generically.

\begin{lemma}[Reachability correctness]\label{lem:reachability}
After \(m\) sweeps, a node \(j\) is marked exactly when there is a directed dependency path from \(t\) to \(j\), allowing the empty path at \(t\).
\end{lemma}

\begin{proof}
For soundness, marks are introduced only at the target or by following one dependency edge from an already marked node.  Induction on the sweep number therefore turns every mark into a path.

For completeness, let \(j\) be at path distance \(d\) from \(t\).  Induction on \(d\) shows that \(j\) is marked after at most \(d\) sweeps.  A simple path has length below \(m\), because the graph is acyclic.  Hence \(m\) sweeps suffice.
\end{proof}

Filtering the marked nodes by boundary kind gives exactly \(R(t)\).  Inputs are not residuals, and derived nodes are expanded to their predecessors.

\subsection{Irrelevance and monotonicity}

\begin{proposition}[Irrelevant boundaries are absent]
If a boundary node \(b\) is not reachable from the target, then \(b\notin R(t)\) and its formula is not needed in the compiled target sequent.
\end{proposition}

\begin{proof}
Immediate from the recursive definition of \(R\), or from \cref{lem:reachability}.
\end{proof}

\begin{proposition}[Monotonic progress]
Replacing a boundary node by a derived node whose predecessors have residual set \(S\) replaces that boundary in every downstream residual by \(S\), leaving all local proofs downstream unchanged.
\end{proposition}

\begin{proof}
The recursive equation for the improved node changes from \(R(b)=\{b\}\) to \(R(b)=S\).  Every downstream residual is obtained by finite union, so substitution propagates mechanically.  The formulas and local derivations of downstream nodes are unchanged; only the context from which the same composed proof is read changes.
\end{proof}

This is the formal sense in which work on producers is modular.

\subsection{Declared boundary implications}

Sometimes a library knows that one boundary implies another.  For example, a complete finite list implies a height bound by taking a maximum, or a stronger modulus implies a weaker one.  Such knowledge should not be hard-coded into residual extraction.  Instead, add a derived node carrying the embedded implication proof.  Dependency expansion will then select the stronger primitive assumption and eliminate the derived one.

This design keeps normalization auditable.  Every nontrivial reduction between boundaries is itself a theorem, not a meta-level heuristic.

\subsection{Alternative producers}

A target may have two independent proof routes.  A plain DAG whose node requires all predecessors cannot represent disjunctive availability directly.  \CIFER\ provides two safe encodings.

\paragraph{Finite route selection.}
Create one derived target for each route, then a final node whose local proof uses an explicit route certificate \(r\) and bounded case analysis.  The producer boundary says that one route certificate and its evidence are available.  This is appropriate when the route is decidable from finite input data.

\paragraph{Relational route cloud.}
Return all successful routes as a finite relation and apply a stable observer to their common conclusion.  This is appropriate when selecting a canonical route is unnecessary or unstable.

What is not permitted is an implicit meta-level choice of whichever unproved route happens to be convenient.

\subsection{Cyclic theorem dependencies}

The proof graph itself must remain acyclic.  A logical cycle
\[
 A\text{ from }B,\qquad B\text{ from }A
\]
contains no proof unless one direction has an independent base, an induction principle, or a least-fixed-point theorem already represented in the object theory.  Such cycles must be packaged as one trace node with an invariant or resolved by an explicit induction.  Only the demand graph may use the finite fixed-point construction of \cref{sec:demand-closure}.

\subsection{Producer/verifier factorization theorem}

Residual compilation makes the following common pattern precise.

\begin{theorem}[Producer--verifier factorization]\label{thm:producer-verifier}
Suppose a verifier graph compiles to
\[
 B(x,e),\Gamma(x)\vdash G(x,e),
\]
and a producer has an embedded proof
\[
 \Gamma(x)\vdash B(x,p(x))
\]
for a primitive-recursive term \(p\).  Then
\[
 \Gamma(x)\vdash G(x,p(x)).
\]
\end{theorem}

\begin{proof}
Instantiate the verifier theorem at \(e=p(x)\), then cut in the producer proof.
\end{proof}

The theorem is elementary, but it captures a crucial engineering boundary.  A difficult producer can be developed and audited independently from a large downstream verification library.

\subsection{Residuals as research questions}

A residual formula should be stated at the strongest reusable level justified by the consumer.  For example, a Diophantine consumer may not need ``the general effective Faltings theorem.''  It may need:
\begin{quote}
Given the code of this fixed curve family and parameter \(u\), return a natural number \(H(u)\) and prove that every rational point has canonical numerator and denominator height at most \(H(u)\).
\end{quote}
This formula is narrower, computationally explicit, and directly connectable to finite enumeration.  The residual report thus becomes a list of concrete research questions with exact downstream consequences.

\subsection{Implementation invariant}

At every prefix of a topological fold, maintain:
\[
 \forall A\in\textit{done},\quad \mathcal B_t,\Gamma\vdash A,
\]
where \(\mathcal B_t\) is the already computed target residual.  A new node's local derivation has context \(\textit{deps}\mathbin{++}(\mathcal B_t\mathbin{++}\Gamma)\).  Simultaneous cut replaces the dependency hypotheses with stored derivations.  This is the invariant implemented by \code{ResidualBuild} in the accompanying Lean outline.

One could instead store the smaller node-specific residual \(\mathcal B_i\) for each node.  That representation is more precise but requires context transport when composing nodes.  The target-residual schedule is simpler: every reachable node is proved under one common context, and weakening supplies unused boundaries.  The paper's mathematical theorem justifies both implementations.

\section{Boolean reflection, bounded search, and finite enumeration}

\subsection{Boolean conventions}

Use \(0\) and \(1=S0\) as false and true.  Fix primitive-recursive operations
\[
 \operatorname{not},\quad \operatorname{and},\quad \operatorname{or},\quad
 \operatorname{imp},\quad \operatorname{eqb},\quad \operatorname{leb}
\]
with truth-table equations.  For example,
\[
 \operatorname{and}(a,b)=
 \begin{cases}
 1,&a=1\text{ and }b=1,\\
 0,&\text{otherwise}.
 \end{cases}
\]
The actual \iPRA\ development need not use case-defined notation; every operation is expressed by primitive recursion or a selector term.  The important facts are equations showing that outputs are Boolean and reflect the intended decidable relation.

\subsection{Equality and order reflection}

For natural numbers, define truncated subtraction and zero test primitively recursively.  Then
\[
 \operatorname{eqb}(a,b)=1
 \quad\Longleftrightarrow\quad a=b
\]
and
\[
 \operatorname{leb}(a,b)=1
 \quad\Longleftrightarrow\quad a\le b.
\]
Both directions are provable by quantifier-free induction.  The repository already contains arithmetic comparison modules that should be reused rather than rebuilt in each application.

For signed integers and rationals, reflection is reduced to natural arithmetic on canonical codes.  Rational equality
\[
 \frac{a}{b}=\frac{c}{d}
\]
with positive denominators is checked by signed cross multiplication \(ad=cb\).  Rational order is similarly checked without division.  Canonical normalization is useful for output but not required for equality reflection if a certified setoid is used.

\subsection{Fixed-formula Booleanization}

Let \(A(\vec x)\) be a fixed quantifier-free formula.  Construct a Boolean term \(\beta_A(\vec x)\) recursively:
\begin{align*}
 \beta_{s=t} &= \operatorname{eqb}(s,t),\\
 \beta_{\bot} &= 0,\\
 \beta_{B\wedge C} &= \operatorname{and}(\beta_B,\beta_C),\\
 \beta_{B\vee C} &= \operatorname{or}(\beta_B,\beta_C),\\
 \beta_{B\to C} &= \operatorname{imp}(\beta_B,\beta_C).
\end{align*}

\begin{theorem}[Fixed quantifier-free reflection]\label{thm:qf-reflection}
For every fixed quantifier-free formula \(A\), one can effectively construct a primitive-recursive Boolean term \(\beta_A\) and embedded derivations
\begin{align*}
 A &\to \beta_A=1,\\
 \beta_A=1 &\to A.
\end{align*}
Moreover \(\iPRA\) proves \(\beta_A=0\vee\beta_A=1\).
\end{theorem}

\begin{proof}
Recurse at the host level over the fixed syntax tree of \(A\).  Atomic equality is the equality-reflection theorem.  The false formula is immediate from successor nonzeroness.  For conjunction and disjunction, use the Boolean truth tables and the induction hypotheses.  For implication, reason intuitionistically as follows.

If \(B\to C\), inspect the Boolean value of \(\beta_B\).  If it is \(0\), then Boolean implication is \(1\).  If it is \(1\), reflection gives \(B\), hence \(C\), and then \(\beta_C=1\); again Boolean implication is \(1\).  Conversely, assume \(\operatorname{imp}(\beta_B,\beta_C)=1\) and assume \(B\).  Reflection gives \(\beta_B=1\).  The truth table then yields \(\beta_C=1\), and reflection gives \(C\).  Thus \(B\to C\).

Booleanness is proved simultaneously from the truth-table equations.  All syntax recursion occurs in the compiler; the generated proof contains only the fixed cases of \(A\).
\end{proof}

The implication case illustrates why intuitionistic logic is sufficient.  Decidability belongs to the Boolean value of the fixed quantifier-free formula, not to arbitrary formulas.

\subsection{Strictly bounded formulas}

Extend the grammar with
\[
 \exists z<t(\vec x)\,B(z,\vec x),
 \qquad
 \forall z<t(\vec x)\,B(z,\vec x),
\]
where \(B\) is already Booleanized.  Define finite folds
\begin{align*}
 E(0)&=0,&E(k+1)&=\operatorname{or}(E(k),\beta_B(k)),\\
 U(0)&=1,&U(k+1)&=\operatorname{and}(U(k),\beta_B(k)).
\end{align*}
Then use \(E(t)\) and \(U(t)\) as the Booleanizers.

\begin{theorem}[Bounded reflection]\label{thm:bounded-reflection}
Every fixed strictly bounded formula has a primitive-recursive Booleanizer with the same two reflection directions as in \cref{thm:qf-reflection}.
\end{theorem}

\begin{proof}
For the existential fold, prove by quantifier-free induction on \(k\) that
\[
 E(k)=1\quad\Longleftrightarrow\quad \exists z<k\;B(z).
\]
The induction step uses the disjunction
\[
 (\exists z<k\;B(z))\vee B(k).
\]
The bounded existential on the right is a human abbreviation for the finite fold already built at stage \(k\); the object invariant may be stated entirely through the Boolean equation \(E(k)=1\) plus the least-witness state described below.  Alternatively, prove both directions by iterating local introduction and elimination proofs along the fixed fold.

The universal case is analogous with conjunction.  Nested bounded quantifiers are handled by host recursion on formula structure; every inner Booleanizer is fixed before the outer fold is generated.
\end{proof}

\subsection{Least bounded search state}

Let \(P(z,\vec x)\) be represented by a Boolean checker \(p(z,\vec x)\).  Search the interval \([0,b)\).  A convenient state is a pair \((f,w)\):
\begin{itemize}
\item \(f=0\) means no witness has been seen;
\item \(f=1\) means \(w\) is the least witness seen so far.
\end{itemize}
Initialize \((0,0)\).  At stage \(k\), if \(f=1\), preserve the state.  If \(f=0\) and \(p(k)=1\), set \((1,k)\).  Otherwise preserve \((0,0)\).

All branches are primitive recursive because \(p\) is Boolean.  The invariant at stage \(k\) is the conjunction of the following quantifier-free or Booleanized clauses:
\begin{enumerate}
\item \(f=0\vee f=1\);
\item if \(f=0\), then \(p(j)=0\) for every \(j<k\);
\item if \(f=1\), then \(w<k\), \(p(w)=1\), and \(p(j)=0\) for every \(j<w\).
\end{enumerate}
The bounded universal clauses are represented by accumulated Boolean conjunctions, so the actual invariant is quantifier-free.

\begin{theorem}[Least bounded search]\label{thm:least-search}
The search state after \(b\) steps satisfies exactly one of:
\begin{enumerate}
\item \(f=0\) and no \(z<b\) satisfies \(P(z)\);
\item \(f=1\), \(w<b\), \(P(w)\), and no \(z<w\) satisfies \(P(z)\).
\end{enumerate}
All conclusions have embedded \iPRA\ derivations generated from one quantifier-free induction invariant.
\end{theorem}

\begin{proof}
The base state satisfies the invariant vacuously.  For the step, split on \(f\).  If a witness was already found, preservation is immediate.  Otherwise split on \(p(k)\).  If it is false, extend the accumulated no-witness conjunction.  If it is true, store \(k\); all earlier checks are false by the old invariant, so minimality holds.  Boolean exhaustiveness supplies the finite cases intuitionistically.

Apply quantifier-free induction to obtain the invariant at \(b\).  Reflect the final Boolean clauses to the stated formulas.
\end{proof}

\subsection{Complete bounded enumeration}

To enumerate all witnesses below \(b\), scan the interval and append \(k\) exactly when \(p(k)=1\).  Encode the output list in a bounded sequence state.  Maintain:
\begin{enumerate}
\item every stored entry is below the current stage and satisfies \(P\);
\item entries are strictly increasing;
\item every \(j<k\) satisfying \(P\) occurs exactly once;
\item the list length is at most \(k\).
\end{enumerate}
Each clause is a bounded Boolean relation on the coded state.

\begin{theorem}[Bounded enumeration]\label{thm:bounded-enumeration}
For every fixed decidable predicate \(P\), there is a primitive-recursive function \(L(b,\vec x)\) such that \iPRA\ proves:
\begin{align*}
 &\forall i<|L|\;\bigl(L_i<b\wedge P(L_i)\bigr),\\
 &\forall z<b\;\bigl(P(z)\to\exists i<|L|\;L_i=z\bigr),\\
 &\text{the entries of }L\text{ are strictly increasing.}
\end{align*}
\end{theorem}

\begin{proof}
Use the scan and invariant above.  Appending at the current stage preserves increasing order.  Soundness of old entries is preserved; soundness of the new entry follows from the branch condition.  Completeness below \(k+1\) splits a witness \(z<k+1\) into \(z<k\) or \(z=k\).  The first case uses the old invariant, and the second uses the append branch.  All comparisons and list accesses are bounded and Booleanized.
\end{proof}

\subsection{Finite maximum and argmin}

Finite extrema are similarly primitive recursive.  For a nonempty list \(L\) and a natural-valued score \(h\), scan with state \((m,i)\) recording the maximum score and its first index.  The invariant states that \(m\) is the maximum of the processed prefix and \(i\) is its least attaining index.  A minimum uses the same pattern after initialization at the first entry.

These routines support:
\begin{itemize}
\item list-to-height conversion by maximum;
\item stable finite-net infima and suprema;
\item canonical choice \emph{within a fixed finite cloud};
\item lexicographically least viable paths;
\item finite optimization certificates.
\end{itemize}
The last two are constructive because the search domain is explicitly finite.  They should not be confused with a global selector over an ideal infinite object.

\subsection{Bounded descent}

Suppose a transition \(T(s)\) strictly decreases a natural rank \(r(s)\) whenever a terminal predicate is false.  Then at most \(r(s_0)\) active steps are possible.  Run exactly that many iterations, freezing the state after termination.  The invariant states:
\begin{enumerate}
\item the state is valid;
\item if terminal, it remains terminal and frozen;
\item if not terminal, the rank has decreased by at least the number of active steps;
\item the public relation between the current state and the input is preserved.
\end{enumerate}
This compiles Euclidean descent, coefficient reduction, finite group descent, lexicographic descent after mixed-radix ranking, and many normalization procedures.

\subsection{Mixed-radix rank for fixed tuples}

For a tuple \((a_0,\ldots,a_{d-1})\) with coordinate caps \(0\le a_i<C_i\), define
\[
 \rho(a)=a_0+C_0a_1+C_0C_1a_2+\cdots+
 \left(\prod_{i<d-1}C_i\right)a_{d-1}.
\]
For fixed \(d\), \(\rho\) is primitive recursive.  Lexicographic decrease with the highest coordinate most significant implies strict decrease of \(\rho\), provided each coordinate remains under its cap.  Thus fixed-dimensional lexicographic descent receives an ordinary natural fuel.  The cap proof belongs in the invariant.

\subsection{Why bounded search solves many apparent logical problems}

Finite mathematical statements are often written with existential quantifiers even when a bound is available.  Once the bound and a decidable verifier are explicit, existential introduction is reduced to least bounded search.  Similarly, a finite universal is a Boolean fold, and a finite disjunction is a bounded route certificate.  This does not trivialize the mathematics: obtaining the correct bound may be the hard residual.  It does make the downstream constructive logic systematic.

\section{Capped monotone demand closure}
\label{sec:demand-closure}

\subsection{Why demand propagation needs its own theorem}

In a simple acyclic numerical program, one computes required input precision by traversing the graph backwards.  Real proof interfaces are often less simple.  A division node needs a lower bound on the denominator; obtaining that lower bound may require a finer approximation; the finer approximation changes the arithmetic size bound; and the size bound influences later precision.  Similar feedback occurs in root isolation, interval Newton methods, finite nets, normalization, and analytic estimates.

This feedback is not a logical proof cycle.  It is a finite system of resource inequalities.  \CIFER\ solves it before forward execution.

\subsection{Finite demand space}

Fix \(d\) demand coordinates and caps \(C_0,\ldots,C_{d-1}\).  The demand space is
\[
 D=\prod_{i<d}\{0,1,\ldots,C_i\}
\]
with componentwise order.  A demand vector may contain requested bits of precision, loop fuel, coefficient height, finite-net depth, denominator exponent, list length, or any other natural resource.

Let \(s\in D\) be the externally requested seed.  Let
\[
 F:D\to\Nat^d
\]
be the aggregate backward transformer.  The intended constraints are \(d\ge s\) and \(d\ge F(d)\).  Define the inflationary operator
\[
 H(d)=d\vee F(d),
\]
where \(\vee\) is componentwise maximum.  A valid closure is a fixed point of \(H\) below the cap.

\subsection{Assumptions}

The finite closure theorem uses three explicit assumptions.

\begin{enumerate}
\item \textbf{Monotonicity:} if \(a\le b\), then \(F(a)\le F(b)\).
\item \textbf{Cap preservation along the trace:} starting at \(s\), every computed \(H^k(s)\) lies below \(C\).
\item \textbf{Primitive recursiveness:} the fixed-dimensional transformer and componentwise maximum are primitive recursive.
\end{enumerate}

Cap preservation may be established in several ways.  A coarse application theorem may prove \(F(C)\le C\), which by monotonicity makes the entire box invariant.  Alternatively, the executable trace can reject cap overflow and return a residual obligation requesting a larger cap.  The framework never silently truncates an overflowing demand, because truncation could invalidate downstream accuracy.

\subsection{Rank argument}

Define
\[
 \rho(d)=\sum_{i<d} d_i.
\]
If \(H(d)\ne d\), then at least one coordinate strictly increases and no coordinate decreases, so
\[
 \rho(H(d))>\rho(d).
\]
Every point below the cap satisfies \(\rho(d)\le\rho(C)\).  Thus there are at most \(\rho(C)-\rho(s)\) strict steps.

\begin{theorem}[Capped stabilization]\label{thm:demand-stabilization}
Assume the iterates \(d_0=s\), \(d_{k+1}=H(d_k)\) remain below \(C\).  Then for some
\[
 k\le \rho(C)-\rho(s)
\]
we have \(d_{k+1}=d_k\).  Running
\[
 T=\rho(C)-\rho(s)+1
\]
steps and freezing after the first fixed point returns a fixed point.
\end{theorem}

\begin{proof}
If every one of the first \(T\) transitions were strict, the rank would increase at least \(T\) times, giving
\[
 \rho(d_T)\ge\rho(s)+T>\rho(C),
\]
contrary to cap preservation.  Therefore some transition is stationary.  Since \(H(d)=d\), all later iterates remain equal to \(d\).
\end{proof}

The proof needs no classical principle and no appeal to an abstract lattice theorem.  It is a bounded natural descent argument suitable for quantifier-free induction.

\subsection{Least fixed point}

\begin{theorem}[Least demand closure]\label{thm:least-demand}
Assume \(F\) is monotone and the capped iteration stabilizes at \(d^*\).  If \(e\in D\) satisfies \(s\le e\) and \(F(e)\le e\), then \(d^*\le e\).  Hence \(d^*\) is the least prefixed point above the seed.
\end{theorem}

\begin{proof}
Induct on \(k\) to show \(d_k\le e\).  The base is \(s\le e\).  If \(d_k\le e\), monotonicity gives \(F(d_k)\le F(e)\le e\), and also \(d_k\le e\).  Therefore
\[
 d_{k+1}=d_k\vee F(d_k)\le e.
\]
At stabilization, \(d^*=d_k\le e\).
\end{proof}

Leastness matters because it prevents gratuitous precision blowup.  It also gives a canonical finite resource assignment without choosing among incomparable solutions.

\subsection{Object-language trace invariant}

For fixed dimension, encode a demand vector as one natural number with primitive-recursive projections and updates.  The closure trace uses state
\[
 (k,d,f),
\]
where \(k\) is the step counter, \(d\) the current vector, and \(f\) a Boolean fixed flag.  The transition computes \(d'=H(d)\); if \(d'=d\), it sets \(f=1\) and freezes; otherwise it stores \(d'\).

A quantifier-free invariant records:
\begin{enumerate}
\item \(d\le C\);
\item \(s\le d\);
\item \(d\) is the declared \(k\)-th frozen iterate;
\item the trace is componentwise inflationary;
\item if \(f=1\), then \(H(d)=d\);
\item if \(f=0\), then \(\rho(d)\ge\rho(s)+k\) for active steps, up to the first fixed point.
\end{enumerate}
All vector comparisons are finite Boolean conjunctions.  Apply \code{PRA.qfInd} for the explicit fuel \(T\), then use the rank contradiction to show the final flag is true.

\subsection{Demand obligations as residuals}

A cap is mathematical data.  If the proof cannot establish that the demand transformer remains below a proposed cap, the correct output is not an arbitrary large hard-coded integer.  There are two honest alternatives:

\begin{enumerate}
\item return a residual boundary formula asserting a sufficient cap and its preservation proof;
\item return a partial conditional result: if closure stays below \(C\), then the downstream theorem holds at the requested precision.
\end{enumerate}

For many fixed finite applications, a crude primitive-recursive cap is easy.  For deep analytic applications, the cap may encode exactly the missing effective constant and should remain visible.

\subsection{A three-coordinate example}

Suppose a consumer needs output precision \(a\), an intermediate normalization needs \(b\), and a supplied denominator certificate has strength \(c\).  The constraints are
\[
 a\ge b+1,
 \qquad
 b\ge c+2,
 \qquad
 c\ge q.
\]
Starting from \((0,0,q)\), use
\[
 F(a,b,c)=(b+1,c+2,q).
\]
The inflationary iteration stabilizes at
\[
 (q+3,q+2,q).
\]
For \(q=3\), the reference implementation obtains \((6,5,3)\).  The example is trivial algebraically, but it models a common pattern: the demanded precision at one node depends on a bound produced downstream.

\subsection{Combining acyclic and cyclic demand propagation}

Most demand edges are acyclic.  It is wasteful to put the whole graph into one large fixed-point system.  Collapse strongly connected components of the demand graph.  The component graph is a DAG.  Process components in reverse topological order:

\begin{itemize}
\item a singleton without a self-loop uses direct backward propagation;
\item a cyclic component uses capped monotone closure on only its coordinates;
\item outgoing demands are then propagated to earlier components.
\end{itemize}

Strongly connected components are computed at the host level on the fixed graph.  The object theory receives the resulting fixed vector and, when desired, the finite closure trace and its verification theorem.

\subsection{Changing-parameter recursion}

Some algorithms change a parameter that affects future precision, such as a shrinking interval width, an evolving degree bound, or a denominator estimate.  The safest compilation is a two-pass bounded trace:

\begin{enumerate}
\item a reverse or summary pass computes the demand attached to each time step;
\item a forward replay runs the actual state transition using those precomputed demands.
\end{enumerate}

Encode both arrays in one finite state if necessary.  The invariant says that the backward table satisfies every local demand constraint and that the forward prefix is valid at the stored demand.  This prevents an informal proof from choosing precision adaptively in a way that is not represented by a primitive-recursive term.

\subsection{Demand closure and proof compilation commute}

Demand closure determines terms and bounds used in local node formulas.  Therefore it must be completed before the final object formulas are frozen.  Once the fixed point is computed, graph specialization emits concrete terms at those demands.  The proof compiler then sees an ordinary acyclic graph.

It is possible to internalize the demand trace itself and derive the fixed-point equations in \iPRA.  This is desirable when the demand assignment appears in the public theorem or when the generator must be independently audited.  It is not necessary to represent the host graph algorithm as a universal object program; only the fixed specialized recurrence is internalized.

\subsection{Failure diagnostics}

The finite closure algorithm has useful explicit failure modes:

\begin{description}[style=nextline]
\item[Nonmonotone transformer.] The least-fixed-point theorem is unavailable.  Split cases or strengthen the state until monotonicity holds.
\item[Cap overflow.] Return the offending coordinate and a residual request for a larger cap or stronger side certificate.
\item[Dimension mismatch.] This is a host graph error, not mathematics.
\item[Unstable after rank fuel.] Either the cap invariant or the implementation is wrong; the rank proof makes this a checkable contradiction.
\end{description}

These diagnostics are preferable to a proof script that fails deep inside a numerical inequality with no indication of the missing resource theorem.

\section{Finite traces and object-level proof compilation}

\subsection{Trace normal form}

A large class of constructive arguments can be put in the following form.  There is a coded state space, an initial state \(s_0(x)\), a primitive-recursive transition \(T(k,s,x)\), a fuel \(N(x)\), a quantifier-free invariant \(I(k,s,x)\), and a final extractor.  Define
\[
 s_{k+1}=T(k,s_k,x).
\]
The proof obligations are:
\begin{align}
 P(x)&\to I(0,s_0(x),x),\label{eq:trace-base}\\
 P(x)\wedge I(k,s_k,x)&\to I(k+1,T(k,s_k,x),x),\label{eq:trace-step}\\
 P(x)\wedge I(N,s_N,x)&\to G(x).\label{eq:trace-finish}
\end{align}
The precondition \(P\) may contain quantifiers.  It remains in the permanent context.  Only \(I\) must be quantifier-free.

\begin{theorem}[Trace compilation]\label{thm:trace-compilation}
Given embedded derivations of \eqref{eq:trace-base}--\eqref{eq:trace-finish}, \iPRA\ derives \(P(x)\to G(x)\).
\end{theorem}

\begin{proof}
Work under hypothesis \(P(x)\).  Apply \code{PRA.qfInd} to the formula \(I\), using the base and step derivations.  Universal elimination at the primitive-recursive term \(N(x)\) gives \(I(N,s_N,x)\).  Cut this into the final extractor \eqref{eq:trace-finish}, then discharge \(P\) by implication introduction.
\end{proof}

The theorem is embodied by \code{CIFER.loop} in the Lean outline.  It is the main route by which executable finite algorithms become object-theory theorems.

\subsection{State design}

The success of a trace proof depends more on state design than on tactic automation.  The state should contain every finite quantity needed to make the induction step local.  Typical components are:

\begin{itemize}
\item the current algorithmic value;
\item counters and bounds;
\item Boolean accumulators for bounded universal facts;
\item the least witness or best candidate seen so far;
\item a rank or remaining fuel;
\item a certificate accumulator, such as Bezout coefficients or a path predecessor table;
\item an error numerator and denominator bound;
\item a frozen terminal flag.
\end{itemize}

A common formalization failure is to keep too little state and then require an unbounded historical statement in the induction step.  The repair is to store a Boolean summary of the history.  For instance, instead of the formula ``every earlier candidate failed,'' store the fold
\[
 a_k=\bigwedge_{j<k}\neg p(j)
\]
and prove its update equation.

\subsection{Freezing after termination}

Many algorithms terminate before their coarse fuel is exhausted.  To retain primitive recursion with a fixed bound, add a Boolean \(done\).  The transition is
\[
 T'(s)=
 \begin{cases}
 s,&done(s)=1,\\
 T(s),&done(s)=0.
 \end{cases}
\]
The invariant includes
\[
 done(s)=1\to T'(s)=s.
\]
Then the state at the full fuel equals the actual terminal state.  This pattern removes partiality and avoids reasoning about the least termination time unless that time is itself requested.

\subsection{Trace reflection versus trace replay}

There are two ways to internalize a computation.

\paragraph{Recompute.}
Use the primitive-recursive state function \(s_k\) directly and prove its invariant.  This is compact when the transition is simple.

\paragraph{Replay supplied evidence.}
Accept a coded trace \(e=(s_0,\ldots,s_N)\), verify every adjacent transition, and derive the endpoint theorem.  This is useful when an external solver or expensive host computation generated the trace.  The verifier is still primitive recursive and the theorem is conditional on its acceptance.

The second pattern is proof-carrying evidence.  It does not trust the external producer.  It can dramatically simplify object arithmetic because checking a step is often easier than reconstructing it.

\subsection{Proof graph compilation}

Let the target residual context \(B\) be fixed.  Maintain a list \(D\) of compiled node formulas and the invariant
\[
 \forall A\in D,\quad B,\Gamma\vdash A.
\]
For a new node \(A\) with dependency sublist \(E\subseteq D\), the local proof has type
\[
 E,B,\Gamma\vdash A.
\]
Simultaneous cut replaces every formula of \(E\) by the stored derivation, while formulas in \(B,\Gamma\) are retained by the hypothesis rule.  The new schedule proves every formula in \(A::D\).

\begin{theorem}[Topological proof fold]\label{thm:proof-fold}
A finite sequence of well-typed local node derivations compiles to a derivation of every node from the common residual and permanent context.
\end{theorem}

\begin{proof}
Induct on the node sequence.  The empty schedule is vacuous.  The step is exactly simultaneous cut as described above.  Old node derivations remain available by the induction hypothesis.
\end{proof}

This theorem is logically simpler than the residual theorem because the common boundary context has already been computed.  Its Lean representation is the structure \code{ResidualBuild}.

\subsection{Local proof constructors}

A node generator should return more than a formula.  A useful host record contains:

\begin{lstlisting}
structure CompiledNode where
  formula    : Fm PTm n
  dependency : List Nat
  localProof : dependencyFormulas ++ boundary ++ Gamma
                 |- [PRAAx] formula
  evidence   : Option Nat
  sourceHash : String
\end{lstlisting}

The \code{sourceHash} is not mathematically necessary, but it permits audit tooling to check that a proof object corresponds to the intended generated source.  Evidence may be omitted when the node is a pure logical or algebraic lemma.

\subsection{Substitution and parametric theorems}

A local derivation is often proved once with variables and then instantiated at concrete terms.  The ConciseIPRA deduction system is closed under simultaneous substitution.  Therefore a library lemma
\[
 \Gamma(\vec x)\vdash A(\vec x)
\]
can be instantiated by a term tuple \(\vec t\) to obtain
\[
 \Gamma(\vec t)\vdash A(\vec t).
\]
This is how generic bounded-search, list, and arithmetic lemmas become application nodes without reproving their logical structure.

Care is required under binders.  The checked syntax library already provides capture-avoiding simultaneous substitution and the needed shift/instantiation laws.  Application code should use those operations rather than implement ad hoc de Bruijn manipulation.

\subsection{Equality transport}

Finite calculations often produce a term definitionally different from the term expected by a theorem.  The deduction system includes a simultaneous Leibniz rule: pointwise provably equal substitutions may be exchanged in any formula.  Derived helpers provide symmetry, transitivity, and congruence.

This supports a useful separation:

\begin{itemize}
\item arithmetic nodes prove equations between concrete terms;
\item logical nodes transport formulas along those equations;
\item the final theorem need not expose the implementation term.
\end{itemize}

The separation also helps proof compression.  A normalized term or certified rewrite system can reduce the number of equality transports without changing the interface theorem.

\subsection{One large trace versus many local traces}

PFIG-style normalization suggests combining a query-specific graph into one primitive recursion with one large invariant.  Local proof graphs suggest compiling each bounded loop separately.  Both are valid, and the choice is pragmatic.

\paragraph{One large trace} can reduce repeated coding and make demand replay explicit.  It may, however, produce a very large invariant and make local debugging difficult.

\paragraph{Many local traces} provide reusable lemmas and smaller invariants.  Their composition introduces more cut steps and intermediate formulas.

\CIFER\ permits a hybrid.  Strongly coupled numerical nodes are fused into a trace component; semantically distinct components remain separate proof nodes.  The residual and demand passes are independent of this fusion choice.

\subsection{Proof-term size and checker simplicity}

A direct natural-deduction proof graph may repeat large subderivations after substitution.  Three techniques control size without changing trust assumptions:

\begin{enumerate}
\item use named Lean definitions so the kernel shares expression structure;
\item package recurring object derivations as lemmas and instantiate them;
\item normalize arithmetic certificates externally but verify compact identities internally.
\end{enumerate}

The trusted checker remains Lean's kernel plus the existing embedded calculus definitions.  The host compiler can be large and optimized because its output is independently checked.

\subsection{Axiom and endpoint audit}

A successful build should automatically report:

\begin{enumerate}
\item the exact type of the public theorem;
\item the transitive Lean axioms reported by \code{\#print axioms};
\item occurrences of \code{sorry}, \code{axiom}, and \code{unsafe} in new source;
\item the residual boundary list;
\item the hash of generated graph metadata;
\item the toolchain and upstream commit.
\end{enumerate}

A clean native theorem with no \code{sorry} is not enough if its type is not \code{Prf PRAAx}.  Conversely, ordinary Lean axioms such as quotient soundness in the metatheory do not become object-theory axioms merely because the embedding is implemented in Lean; the audit should report both layers clearly.

\subsection{The final export}

Once the compiler produces
\[
 \Gamma\vdash_{\PRAAx}G,
\]
the checked theorem \code{pra\_iff\_native} gives a native derivation of the translated sequent.  Round-trip theorems show that translation preserves the intended proof content.  This final step is small by design.  All mathematical work has already occurred in the ordinary derivation.

\section{Finite relational witnesses and stable observations}

\subsection{Why relations are sometimes the constructive object}

A classical proof often treats an object as uniquely determined up to an equivalence that is harmless for the final theorem.  At finite precision, however, several candidates may be equally valid.  Forcing one canonical candidate can create a discontinuous selector that is stronger than the theorem actually needs.

The constructive response is not to abandon computation.  It is to compute a finite relation of candidates and preserve enough compatibility information for the consumer.  The relation is finite, decidable, and amenable to bounded search.

\subsection{Finite layered systems}

Fix a horizon \(h\).  For each level \(k\le h\), let \(L_k\) be a nonempty finite list of candidates.  Let
\[
 E_k(x,y)
\]
be a decidable compatibility relation between \(x\in L_k\) and \(y\in L_{k+1}\).  A path is a tuple
\[
 x_0,\ldots,x_h
\]
with \(x_k\in L_k\) and \(E_k(x_k,x_{k+1})\) for every \(k<h\).

This schema covers:
\begin{itemize}
\item nested rational enclosures;
\item root clusters across increasing precision;
\item approximate eigenvectors or invariant subspaces;
\item local chart choices;
\item finite covers and refinements;
\item candidates in a bounded descent;
\item proof routes or auxiliary witnesses.
\end{itemize}

\subsection{Backward viability}

Define viable sets backwards:
\begin{align*}
 V_h&=L_h,\\
 V_k&=\{x\in L_k:\exists y\in V_{k+1}\ E_k(x,y)\}.
\end{align*}
Every set is finite and its characteristic table is primitive recursive.

\begin{theorem}[Viability/path equivalence]\label{thm:viability-path}
For \(x\in L_k\), the following are equivalent:
\begin{enumerate}
\item \(x\in V_k\);
\item there is a compatible path \(x=x_k,x_{k+1},\ldots,x_h\) to the horizon.
\end{enumerate}
\end{theorem}

\begin{proof}
Use reverse induction on \(k\).  At \(k=h\), both statements reduce to \(x\in L_h\).  For \(k<h\), membership in \(V_k\) gives a compatible successor \(y\in V_{k+1}\).  Apply the induction hypothesis to extend \(y\) to a path.  Conversely, the first successor of any path belongs to \(V_{k+1}\) by the induction hypothesis, so the first vertex belongs to \(V_k\).
\end{proof}

The reverse induction is compiled as an ordinary forward primitive recursion on reversed time.  The invariant is a finite Boolean equivalence table.

\subsection{Least finite path}

If \(V_0\) is nonempty, choose its least coded element.  Given the chosen \(x_k\), search the finite list \(V_{k+1}\) for the least compatible successor.  Viability proves that this search succeeds.

\begin{corollary}[Canonical finite path]\label{cor:least-finite-path}
Whenever a path to the fixed horizon exists, a primitive-recursive bounded search returns the lexicographically least such path.
\end{corollary}

This is a legitimate canonical choice because the search domain is finite and explicitly bounded.  It does not define a canonical branch through an infinite tree.

\subsection{Stable exact observations}

Let \(O\) be a finite coded observation and \(\omega(x)\) an observer on candidates.  Suppose the verifier proves
\[
 \forall x,y\in V_k\quad \omega(x)=\omega(y).
\]
Then any viable candidate determines the same observation.

\begin{theorem}[Stable-observer elimination]\label{thm:stable-observer}
If \(V_k\) is nonempty and \(\omega\) is constant on \(V_k\), there is a primitive-recursive output \(o\) and an \iPRA\ proof that \(o=\omega(x)\) for every viable \(x\).  No selector coherent across horizons is required.
\end{theorem}

\begin{proof}
Use least bounded search to choose the first viable element \(x_0\) and set \(o=\omega(x_0)\).  For an arbitrary viable \(x\), apply the stability theorem to \(x_0,x\).  The choice is internal to one fixed finite cloud; the conclusion is independent of it.
\end{proof}

Typical stable observations include a multiplicity, a sign separated from zero, a symmetric polynomial of roots, an induced map on a quotient, a finite truth value, and an interval enclosure common to all branches.

\subsection{Approximate observations by diameter}

Exact constancy is often stronger than needed.  Suppose observations are rational codes in a metric and the verifier proves
\[
 \diam\{\omega(x):x\in V_k\}\le 2^{-q}.
\]
Then the least viable observation is a valid \(q\)-approximation to every branch.

\begin{theorem}[Diameter elimination]\label{thm:diameter-elimination}
A nonempty finite cloud with observation diameter at most \(\varepsilon\) determines an \(\varepsilon\)-stable observable by bounded least choice.
\end{theorem}

\begin{proof}
Choose any member \(x_0\) by bounded search.  The diameter inequality directly gives
\[
 d(\omega(x_0),\omega(x))\le\varepsilon
\]
for every member \(x\).  All maxima and pairwise comparisons are finite.
\end{proof}

This theorem is particularly useful for finite nets and root clusters.  The exact internal label can remain ambiguous while the requested numerical observable is determined to the desired precision.

\subsection{Gap-guarded exact collapse}

Suppose candidates carry labels from a finite set \(J\), and different labels are known to be separated by more than \(2\varepsilon\) in observation space.  If a cloud has diameter at most \(\varepsilon\), it cannot contain two different labels.

\begin{theorem}[Exact recovery from a gap]\label{thm:gap-collapse}
Let \(g>0\) be a certified lower separation between observations of distinct labels.  If every valid candidate lies within radius \(e\) of the true observation and \(2e<g\), then all valid candidates have the same label, which can be recovered by bounded search.
\end{theorem}

\begin{proof}
Assume two valid candidates have different labels.  Their observations are at least \(g\) apart.  By the triangle inequality through the true observation, they are at most \(2e\) apart, contradicting \(2e<g\).  Stability then permits exact bounded choice.
\end{proof}

The gap is not a technical nuisance.  It is the positive information that makes a discontinuous label locally stable.

\subsection{Changing horizons without infinite choice}

For each requested precision \(q\), choose a finite horizon \(h(q)\), compute the viability table to that horizon, and let \(a_q\) be the least viable candidate at a designated level \(k(q)\).  The choices for different \(q\) need not lie on one infinite branch.  Nevertheless they can form a coherent name under quantitative conditions.

Assume:
\begin{enumerate}
\item every full-horizon viable cloud at level \(k(q)\) has diameter at most \(\delta(q)\);
\item compatible edges between levels \(r\) and \(r+1\) move observations by at most \(\eta(r)\);
\item the tail \(\sum_{r\ge k(q)}\eta(r)\) has an explicit primitive-recursive bound \(\tau(q)\);
\item \(\delta(q)+2\tau(q)\le 2^{-q}\).
\end{enumerate}

\begin{theorem}[Shrinking finite-horizon coherence]\label{thm:shrinking-coherence}
Under the four assumptions, the sequence of least finite-horizon observations \(\omega(a_q)\) is Cauchy with the displayed modulus, even though no coherent infinite branch has been selected.
\end{theorem}

\begin{proof}
Take precisions \(p,q\) and a common horizon extending both finite computations.  Extend each selected candidate to that common horizon using its viability path.  At a common comparison level, the two extended candidates lie in the same full-horizon viable cloud, so their observations differ by at most \(\delta(r)\).  Moving from the original selected levels to the comparison level costs at most one motion tail for each path.  The triangle inequality gives
\[
 d(\omega(a_p),\omega(a_q))
 \le \tau(p)+\delta(r)+\tau(q).
\]
Choose the level and horizon functions so the right side is bounded by the requested dyadic error.  All paths are finite and chosen by bounded search.
\end{proof}

This theorem is a constructive substitute for an appeal to compactness plus choice of an infinite branch.  It is stronger than arbitrary finite satisfiability because it uses explicit shrinking and motion moduli.

\subsection{Failure of naive least-at-each-level selection}

Choosing the least candidate in each raw level \(L_k\) can fail in two ways.  The candidate may be a dead end, and least candidates at successive levels may be incompatible.  Backward viability fixes the first problem by removing dead ends.  It does not by itself fix changing-horizon jumps.  Shrinking diameter and motion bounds fix the second problem at the observational level.

Thus there are three distinct least-choice operations:
\begin{enumerate}
\item least element of a finite cloud;
\item least path to a fixed finite horizon;
\item least viable candidate in a family of changing finite horizons.
\end{enumerate}
Only the third requires a coherence theorem, and its conclusion is usually observational rather than literal code convergence.

\subsection{Finite auxiliary exhaustion}

A target object may have several possible auxiliary witnesses, none canonically chosen.  Suppose there is an explicit finite list \(A_0,\ldots,A_{r-1}\) such that every target admits at least one auxiliary in the list.  Instead of defining a selector from targets to auxiliaries, run the downstream verifier for every \(A_i\), union the resulting finite candidate sets, and filter by the target relation.

\begin{proposition}[Auxiliary exhaustion]
If every valid target has some auxiliary in a complete finite list and the verifier is sound for each auxiliary, then the union-and-filter construction is sound and complete for targets.
\end{proposition}

\begin{proof}
Soundness follows from the final target filter.  For completeness, take a target and use the finite-cover theorem to obtain an index \(i<r\).  Completeness of the verifier for \(A_i\) puts the target into that branch's candidate set, hence into the union, and the target filter retains it.
\end{proof}

This method is frequently preferable to a global target-to-auxiliary choice function.

\subsection{Stable observations in proof graphs}

A stable observation node has dependencies consisting of:
\begin{itemize}
\item one or more independently produced certificate formulas;
\item a local comparison theorem showing they determine the same observation;
\item optionally a nonemptiness or gap certificate.
\end{itemize}
The local proof may conclude an existential, equality, or approximation formula.  In the object compiler it is an ordinary derived node.  The special terminology records the constructive reason no selector is needed.

\subsection{When relational output is the final theorem}

Not every relation should be eliminated.  A theorem such as ``the roots at precision \(q\) are covered by these disjoint or overlapping rational boxes'' is already an exact finite mathematical statement.  For singular inputs, it may be strictly more informative and more stable than an arbitrary ordering of roots.  Similarly, a finite atlas, a set of compatible local solutions, or a list of possible group representatives can be a satisfactory final output.

The framework therefore treats finite relations as first-class results, not merely intermediate failures to choose.

\section{Reusable mathematical representations}

\subsection{Natural numbers, integers, and rationals}

Natural numbers are the ambient domain.  Signed integers can be represented by pairs \((p,n)\) denoting \(p-n\), with equality by cross addition.  A normalized sign-magnitude code is smaller for output, but pair-difference coding often simplifies primitive recursion because addition and multiplication are polynomial operations on naturals.

A rational code consists of a signed numerator and positive denominator.  Validity checks positivity of the denominator.  Equality and order use cross multiplication.  Arithmetic operations return valid but not necessarily reduced codes; reduction by the Euclidean algorithm is a separate bounded trace when canonical output is required.  This separation avoids inserting gcd normalization into every rational operation.

For a fixed proof, all denominator growth is bounded by primitive-recursive expressions in the input codes and loop fuel.  The demand pass records the bit or height bounds needed by finite inequalities.

\subsection{Finite-dimensional algebra}

Vectors and matrices of fixed dimensions are nested tuples.  Every coordinate operation is primitive recursive.  Matrix multiplication uses a fixed finite sum, determinant uses a fixed permutation sum or elimination trace, and polynomial evaluation uses Horner recursion.

This representation is especially effective for theorem families with externally fixed dimension: Euclidean and projective geometry, small matrix identities, finite Markov chains, finite group tables, and fixed-degree polynomial algebra.  There is no need for an internal type universe of dimensions.  The host compiler specializes the dimension and generates the finite formulas.

For variable but bounded dimension, encode arrays and use bounded loops.  The invariant must include array length and index safety.  A determinant can then be verified from a supplied permutation expansion or elimination certificate; production and verification remain separate.

\subsection{Polynomials and formal power series}

A polynomial of degree at most \(d\) is a coefficient list of length \(d+1\).  Addition, multiplication, derivative, evaluation, and coefficient extraction are bounded operations.  Identities are proved coefficientwise up to an explicit degree cap.  This is often simpler than reasoning extensionally about polynomial functions.

A formal power series theorem usually consumes only finitely many coefficients.  Represent a series by a primitive-recursive coefficient function or a finite prefix producer.  At query order \(N\), all algebra is performed modulo \(X^{N+1}\).  Composition and inversion require explicit constant-term side conditions; Lagrange inversion and recurrence arguments become finite coefficient identities.  No completed infinite series object is needed unless the theorem asks for one.

\subsection{Represented real numbers}

A regular rational name for a real is a primitive-recursive function \(a(q)\) satisfying a fixed Cauchy condition such as
\[
 |a(p)-a(q)|\le 2^{-p}+2^{-q}.
\]
In the object language, one does not quantify over arbitrary functions.  A theorem is a schema for each fixed \code{PR} symbol representing the name, together with embedded proofs of its regularity equations or inequalities.

An interval name uses rational endpoints \(l(q)\le u(q)\), nested or approximately nested, with width at most \(2^{-q}\).  Interval names are often better for inequalities, root enclosures, and branch stability.  Operations carry demand transformers and interval correctness theorems.

Two names represent the same real when their approximants become close at an explicit common precision.  Consumers should prove extensionality: replacing a name by an equivalent one changes the output by no more than the requested tolerance or leaves a finite observation unchanged.

\subsection{Continuous functions on compact domains}

A constructive finite-use interface for \(f:[a,b]\to\mathbb R\) contains:
\begin{enumerate}
\item a pointwise approximation producer for rational inputs;
\item a modulus of uniform continuity on the declared compact interval;
\item optional range or Lipschitz bounds;
\item extensionality certificates.
\end{enumerate}
To compute a uniform bound or approximate extremum at precision \(q\), choose a rational mesh from the modulus, evaluate a finite net, and take a finite maximum or minimum.  The proof combines the net extremum with the continuity error.

The modulus is often the decisive producer.  If it is unavailable, the finite-net theorem remains valid conditionally with the modulus formula in the residual context.

\subsection{Integration and finite sums}

Riemann-style integration on a compact interval uses a finite partition whose mesh is selected from a modulus.  The evidence is a rational sum and an error bound.  Algebraic identities between integrals are proved first for finite sums and then transported through the error estimates.

Probability distributions with finite support are even simpler: expectations, variances, transforms, and convolutions are finite rational sums.  Limit theorems require explicit quantitative approximation statements; the finite distribution algebra itself is fully primitive recursive.

\subsection{Roots and algebraic numbers}

A robust root interface returns a finite list of rational boxes or discs together with:
\begin{itemize}
\item a verifier that each box contains the claimed number of roots, expressed through an appropriate finite certificate;
\item a cover theorem saying all relevant roots lie in the union;
\item optional disjointness or separation data;
\item coefficient and evaluation bounds.
\end{itemize}
At multiple-root collisions, labels need not be stable.  Symmetric functions of roots, total multiplicities, cluster sums, and polynomial remainders are stable observables.  An exact ordered root can be recovered only under a separation certificate.

Algebraic numbers may be represented by a polynomial plus an isolating interval.  Equality is a certified relation, not literal code equality.  Arithmetic operations produce a resultant polynomial and a finite root cloud; a gap or isolating certificate selects the desired component when possible.

\subsection{Linear algebra and spectral data}

For rational matrices in fixed dimension, row reduction, rank, determinant, kernel bases, and characteristic polynomials are finite.  For represented real matrices, use interval entries and finite precision.

Eigenvectors are collision-prone.  Better interfaces include:
\begin{itemize}
\item the characteristic polynomial;
\item clusters of approximate eigenvalues;
\item invariant-subspace projectors under a spectral gap;
\item symmetric spectral observables such as trace, determinant, or cluster sums;
\item finite clouds of normalized approximate eigenvectors.
\end{itemize}
A unique normalized eigenvector requires a simple eigenvalue and a phase/sign convention separated from zero.  Those are gap boundaries, not automatic consequences of approximate matrix data.

\subsection{Geometry}

Rational or algebraic coordinate geometry is particularly well suited to \iPRA.  Incidence, orientation, determinant identities, squared distances, barycentric coordinates, and projective relations reduce to polynomial equalities and inequalities.  Degenerate cases should be stated explicitly as side formulas.

Trigonometric geometry can often be recast algebraically through dot products, Gram matrices, complex unit pairs, or finite polynomial approximations with error bounds.  An angle is not always the best constructive primitive.  For rotations, unit quaternions or orthogonal matrices avoid the unstable exact axis at the identity.

A geometric construction that uses intersections should return a finite relation when parallelism, tangency, or multiplicity can occur.  Gap conditions recover the classical unique point in nondegenerate cases.

\subsection{Finite combinatorics and graph theory}

Finite graphs, hypergraphs, simplicial complexes, automata, and game trees are direct bounded structures.  Existence claims become bounded search with certificate verification.  Dynamic programs and exhaustive searches compile through trace invariants.  Symmetry reduction may be performed by a host canonicalizer, provided the object proof verifies the canonical representative or the final certificate independently.

For large finite computations, a supplied proof certificate is often preferable to internal search.  Examples include a coloring, matching, flow, elimination order, unsatisfiability trace, or decomposition tree.  The object theorem is then ``accepted certificate implies property.''

\subsection{Algebraic structures by finite laws}

A finite group, ring, field, category, or module is represented by operation tables and Boolean law checkers.  Once associativity, identity, inverses, and distributive laws are verified over the finite carrier, all elementwise reasoning is bounded.  A theorem about a specific finite structure can therefore be closed by exhaustive verification plus an embedded reflection proof.

For infinite structures generated by finite data, use a presentation plus a certified normalizer or relational equality.  A global decision procedure is not assumed unless its termination and correctness are proved by a trace.

\subsection{Topology and homology through finite carriers}

Many topological applications in the repository are best understood as finite-combinatorial theorems plus a representation theorem.  A finite simplicial complex, chain group in fixed degree, boundary matrix, and homology certificate are primitive recursive.  Subdivision and simplicial approximation are bounded traces when a subdivision depth is supplied.

The difficult producer may be a finite cover, Lebesgue number, triangulation, or approximation depth.  Once supplied, chain-level calculations and local-to-global gluing can be compiled.  Homology classes use certified setoid equality witnessed by boundaries.

\subsection{Local-to-global interfaces}

A local-to-global theorem starts with a finite cover \(U_0,\ldots,U_{r-1}\), local evidence \(e_i\), and overlap certificates \(c_{ij}\).  A finite gluing verifier checks:
\begin{enumerate}
\item each local object is valid;
\item overlap restrictions agree under the certified equality relation;
\item higher overlap conditions required by the chosen construction hold;
\item the cover is complete for the finite carrier or represented domain at the demanded scale.
\end{enumerate}
The finalizer constructs or verifies a global code and proves its restrictions.  In sheaf-like applications, the exact cover or Lebesgue number may be the residual producer; the gluing algebra is finite.

\subsection{A representation selection rule}

Before formalization, ask what the final consumer observes.  Choose the weakest representation that makes that observation primitive-recursive and extensional.

\begin{itemize}
\item If only a finite truth value is needed, do not construct an exact real witness.
\item If only a symmetric root expression is needed, do not order roots.
\item If only a homology class is needed, retain chain representatives with boundary certificates.
\item If only a numerical bound is needed, use interval enclosures rather than exact radicals.
\item If the theorem needs a complete list, demand a cutoff, not merely a count.
\end{itemize}

This rule is the main practical source of simplification.  It removes unnecessary selectors before they become logical obstacles.

\section{Resolution of recurring application difficulties}

\subsection{Diophantine finiteness: height bounds, lists, and counts}

Let \(R(u,z)\) be a decidable predicate saying that code \(z\) is a rational or integral point on the object with parameter \(u\).  Let \(h(z)\) be a primitive-recursive canonical height on codes, and assume that for each \(H\) the set of valid codes with \(h(z)\le H\) is contained in an explicitly enumerable finite interval.

\begin{theorem}[Height bound to complete list]\label{thm:height-to-list}
If \iPRA\ proves
\[
 \forall z\;(R(u,z)\to h(z)\le H(u)),
\]
then a primitive-recursive bounded filter produces a list \(L(u)\) such that \iPRA\ proves
\[
 \forall z\;\bigl(R(u,z)\leftrightarrow z\in L(u)\bigr),
\]
with membership interpreted through the canonical code convention.
\end{theorem}

\begin{proof}
Use the height bound and coding theorem to compute a finite code cutoff \(C(u,H(u))\).  Enumerate all codes below the cutoff and retain exactly those satisfying the Boolean checker for \(R\).  Soundness follows from filtering.  For completeness, any solution has height at most \(H(u)\), hence has a canonical code below the cutoff, and therefore is inspected and retained.  The bounded enumeration theorem supplies the embedded list proof.
\end{proof}

\begin{theorem}[Complete list to height bound]\label{thm:list-to-height}
If \iPRA\ proves that a finite list \(L(u)\) contains every solution, then the finite maximum of \(h\) on \(L(u)\) is a height bound for every solution.
\end{theorem}

\begin{proof}
Compute \(M=\max\{h(z):z\in L(u)\}\), taking \(M=0\) for an empty list.  Given a solution \(z\), completeness gives an index \(i<|L|\) with \(z=L_i\); the maximum theorem gives \(h(z)\le M\).
\end{proof}

Thus positive height bounds and complete finite lists are constructively interconvertible for a decidable coded point relation.  This resolves the downstream part of the Faltings and Siegel packets.

A cardinality bound
\[
 |\{z:R(u,z)\}|\le N(u)
\]
does not provide either construction.  It says how many solutions may exist, not where to stop searching.  Consequently, \CIFER\ assigns distinct boundary formulas:
\begin{align*}
 \operatorname{CountBound}(u,N),\qquad
 \operatorname{HeightBound}(u,H),\qquad
 \operatorname{CompleteList}(u,L).
\end{align*}
Only the last two are interchangeable by the generic nodes above.

\paragraph{Residual for general Faltings/Siegel use.}
The correct residual is not a vague statement ``make the classical proof effective.''  It is a formula asserting a primitive-recursive height cutoff for the chosen family and canonical coding, with an embedded proof that every point lies below it.  All list production, verification, maximum extraction, and finite consequences then close automatically.

\subsection{Runge and concrete descent subcases}

The repository's constructive Runge and concrete curve packets show that some height producers can be built directly from polynomial identities.  A typical certificate has the form
\[
 D^2f=A^2+R,
 \qquad \deg R<\deg A,
\]

after which an elementary inequality bounds integral solutions.  In \CIFER\:

\begin{enumerate}
\item the polynomial identity is a finite certificate node;
\item coefficient and degree checks are deterministic nodes;
\item the inequality proof is a local arithmetic node;
\item the resulting height bound discharges the generic list residual;
\item bounded enumeration returns the exact solution list.
\end{enumerate}

This decomposition separates a reusable height-to-list compiler from theorem-specific algebra that produces the height.

\subsection{Mordell--Weil descent and rational covering representatives}

A finite descent theorem often has the following data:
\begin{itemize}
\item a coded abelian group operation on rational points;
\item a height \(h\);
\item a finite list \(R_0,\ldots,R_{r-1}\);
\item a covering assertion that every point \(P\) can be written \(P=2Q+R_i\);
\item an inequality \(2h(Q)\le h(P)+C\).
\end{itemize}
Once this data is present, repeated descent reaches a bounded-height region after explicit bit-length fuel.  The finite region can be enumerated; unwinding the trace yields generators and coordinates.

The difficult point is the covering list of actual rational representatives.  Finite Selmer information can bound or describe a quotient abstractly without producing a rational point in every represented class.  \CIFER\ therefore uses a residual formula
\[
 \operatorname{RationalCover}(E,R)
\]

together with a separate verifier for the group equations and descent inequality.  The downstream theorem is
\[
 \operatorname{RationalCover}(E,R),
 \operatorname{DescentBound}(E,R,C)
 \vdash \operatorname{FiniteGenerators}(E,G).
\]
No selector from a point to its cover index is assumed.  Given the universal covering formula, bounded search over the finite list and a bounded rational halving verifier can find the least successful pair \((i,Q)\) at each descent stage.

Concrete curves whose rational points are classified by an independent descent have no residual: the classification supplies the finite list directly.  General curves retain the covering formula visibly.

\subsection{Exact axes, eigenvectors, and branch selectors}

For a proper rotation matrix near the identity, the exact unit axis is unstable: at the identity every axis is valid, and arbitrarily small rotations may have unrelated axes.  A continuous or primitive-recursive exact selector from arbitrary Cauchy names would decide information not constructively available in general.

\CIFER\ offers four statement levels.

\begin{enumerate}
\item \textbf{Quaternion level.}  Return an exact or approximate unit quaternion, allowing its vector part to vanish at the identity.  This avoids selecting an axis when none is distinguished.
\item \textbf{Finite axis cloud.}  At precision \(q\), return finitely many rational axis-angle candidates whose induced rotations approximate the input.
\item \textbf{Approximate stable observable.}  Prove that all viable candidates induce rotations within the requested matrix error; the axes themselves need not converge.
\item \textbf{Gap-conditional exact axis.}  Under a supplied positive lower bound separating the rotation from the identity and other singular cases, normalize the vector part and recover an exact represented axis.
\end{enumerate}

The same pattern applies to a normalized eigenvector at a simple eigenvalue.  The residual is a spectral gap and, if a sign or phase is fixed, a nonzero coordinate bound.  Without the gap, the stable object is the invariant subspace, spectral projector, characteristic polynomial, or finite cloud.

This resolves the formal difficulty by changing the default conclusion, not by pretending a discontinuous selector is primitive recursive.

\subsection{Root ordering and singular branches}

For a polynomial with represented coefficients, roots can collide as coefficients vary.  An exact global ordering of root names is therefore unstable.  The finite relation interface returns rational boxes with multiplicity and coverage certificates.  The consumer can compute:

\begin{itemize}
\item symmetric polynomials;
\item total multiplicity in a region;
\item cluster centers and diameters;
\item signs or labels separated by a gap;
\item resultants and gcd degrees.
\end{itemize}

If a discriminant lower bound or isolating separation is supplied, the gap-collapse theorem yields unique labels.  Otherwise, the cluster relation is the correct exact result.

\subsection{Prime distribution and effective constants}

Finite arithmetic after an explicit analytic estimate is generally compatible with \iPRA.  Given rational constants, cutoffs, and error inequalities, one can compute finite sums, verify divisor identities, clear denominators, and perform bounded prime search.  The difficult producer is often one of:

\begin{itemize}
\item an effective zero-free region;
\item a density estimate with explicit constants;
\item field-dependent discriminant and conductor parameters;
\item control of an exceptional real zero;
\item a mean-value estimate uniform in a conductor range.
\end{itemize}

A \CIFER\ analytic node therefore separates:
\begin{enumerate}
\item \(E(\alpha)\): a boundary formula saying that parameter code \(\alpha\) satisfies the required effective estimate;
\item a finite verifier of every numerical side condition;
\item a local derivation from \(E(\alpha)\) to the desired finite inequality or bounded prime witness.
\end{enumerate}

For constructive proofs such as explicit Dirichlet infinitude or finite sieve inequalities, the estimate producer may itself be elementary and the residual closes.  For Chebotarev, Siegel--Walfisz, or Bombieri--Vinogradov at full requested uniformity, the remaining analytic producer stays explicit.  The framework still completes all denominator clearing, finite identities, and final bounded search.

\subsection{Finite nets and compactness surrogates}

Classical analysis often invokes compactness to choose a finite subcover, a convergent subsequence, or an extremizer.  In a quantitative constructive setting, replace the invocation by explicit finite data:

\begin{itemize}
\item a modulus or total-boundedness function producing an \(\varepsilon\)-net;
\item a finite coverage verifier;
\item a uniform continuity modulus for the consumer;
\item finite extrema or viability on the net.
\end{itemize}

If the total-boundedness modulus is known, approximate extrema and finite subcovers compile directly.  If not, it is a residual.  An ideal compact object is never used as an unexamined source of finite choice.

\subsection{Local-to-global topology}

A typical topological proof has a finite combinatorial core and a difficult approximation or cover theorem.  For example, after a triangulation and a simplicial approximation depth are supplied, chain maps, boundary identities, degrees, and homology calculations are finite integer algebra.  \CIFER\ makes the triangulation/approximation theorem a boundary and compiles the chain-level consequences.

When exact representatives are noncanonical, the final theorem should concern a homology class or induced map, with equality witnessed by a boundary chain.  Stable-observable elimination removes dependence on the chosen simplicial approximation if a finite chain homotopy certificate is supplied.

\subsection{Analysis with singular denominators}

Division by a represented real is not uniformly total unless nonzeroness carries quantitative content.  A proposition \(x\ne0\) alone may not provide a primitive-recursive lower bound on \(|x|\) from an arbitrary name.  The correct interface is a guard
\[
 |x|\ge 2^{-a}
\]
with explicit exponent \(a\).  Then reciprocal precision is a primitive-recursive function of \(q\) and \(a\), and interval arithmetic proves correctness.

Thus denominator guards become resource-bearing evidence.  If a classical theorem assumes only abstract nonzeroness, the constructive reformulation either adds a located gap, weakens the result to a relation, or keeps the gap as a residual.

\subsection{Probability and limit theorems}

Finite probability distributions and moments are straightforward finite sums.  The obstacle in a central limit or concentration theorem is quantitative convergence: a rate, truncation bound, smoothing parameter, or characteristic-function estimate.  \CIFER\ can compile the entire finite Fourier or moment calculation once those parameters are supplied.

The result should be stated at finite precision:
\[
 \left|\Pr(S_n\le x)-\Phi(x)\right|\le 2^{-q}
\]
for an explicit primitive-recursive \(n\ge N(q,\text{data})\), with rational approximants to \(\Phi\).  A purely asymptotic statement without a rate does not by itself give the producer needed by \iPRA.

\subsection{Transcendental functions}

Sine, cosine, exponential, and logarithm can be represented by finite sums with explicit remainder bounds on bounded domains.  Argument reduction, domain guards, and precision propagation are the substantive interfaces.  Rather than formalize a completed real function first, compile the exact polynomial/rational approximation consumed by the theorem.

For trigonometric identities, many statements are better expressed algebraically on a unit-circle pair \((c,s)\) satisfying \(c^2+s^2=1\).  Analytic sine/cosine names then become one producer of such pairs.  This separates polynomial identity proofs from convergence proofs.

\subsection{What the framework has and has not resolved}

The general obstacles divide cleanly.

\begin{longtable}{p{0.32\textwidth}p{0.58\textwidth}}
\toprule
Obstacle & CIFER resolution \\
\midrule
\endhead
Host computation not an iPRA proof & Require local \code{Prf PRAAx} constructors and compose by cut. \\
Hidden precision/fuel & Backward demand plus capped monotone closure. \\
Unbounded finite search & Demand an explicit cutoff; use bounded search and enumeration. \\
Count mistaken for cutoff & Separate count, height, and complete-list formulas. \\
Unstable exact choice & Preserve a finite relation; use stable observables or a gap. \\
Infinite branch choice & Use fixed-horizon viability and shrinking finite-horizon coherence. \\
Universal PR evaluation & Specialize a fixed graph to concrete \code{PR} syntax. \\
Quotient equality & Use certified setoids and explicit equality witnesses. \\
Deep ineffective theorem & Return its exact producer formula as a residual boundary. \\
\bottomrule
\end{longtable}

The last row remains theorem-specific by necessity.  The framework's contribution is that every other row can be solved once, reused, and checked at the actual embedded endpoint.

\section{Worked closed benchmark: the inverse of \(y^3+y\)}

\subsection{Purpose of the benchmark}

A general framework should close at least one nontrivial theorem end to end.  The equation
\begin{equation}\label{eq:cubic-equation}
 y^3+y=x
\end{equation}
is a useful benchmark because it combines signed rational coding, inequalities, a bounded trace, precision demand, a represented-real output, and uniqueness, yet it needs no deep residual producer.

The input is a rational code \(x=a/b\) with \(b>0\).  The output is a primitive-recursive regular name for the unique real solution, together with rational intervals certified to contain it at every requested precision.

\subsection{Strict monotonicity}

Let \(F(t)=t^3+t\).  For rationals \(r<s\),
\begin{align*}
 F(s)-F(r)
 &= (s-r)(r^2+rs+s^2+1).
\end{align*}
The second factor is positive.  Indeed,
\[
 4(r^2+rs+s^2)=(2r+s)^2+3s^2\ge0,
\]
so \(r^2+rs+s^2+1\ge1\).  Hence \(F(s)>F(r)\).

All inequalities can be checked by clearing positive denominators.  The identity is polynomial and can be proved by ring normalization or a short sequence of distributive rewrites in the object arithmetic library.

\begin{lemma}[Cubic monotonicity]\label{lem:cubic-monotone}
For rational codes \(r\le s\), \iPRA\ proves \(F(r)\le F(s)\); if \(r<s\), it proves \(F(r)<F(s)\).
\end{lemma}

The lemma supplies both preservation of the bisection bracket and uniqueness of the represented root.

\subsection{Initial bracket}

Set
\[
 M=|x|+1.
\]
Then
\begin{align*}
 F(M)-x &\ge M-x = |x|+1-x\ge1,\\
 F(-M)-x &\le -M-x =-(|x|+x+1)\le-1.
\end{align*}
Thus
\[
 F(-M)\le x\le F(M).
\]
The interval \(I_0=[-M,M]\) is a certified bracket.  Its width is \(2M\).

\subsection{Exact rational bisection}

Given a bracket \(I_k=[l_k,u_k]\), let
\[
 m_k=\frac{l_k+u_k}{2}.
\]
Use the decidable rational comparison \(F(m_k)\le x\).  If true, set
\[
 I_{k+1}=[m_k,u_k];
\]
otherwise set
\[
 I_{k+1}=[l_k,m_k].
\]
Every operation is primitive recursive on rational codes.  The interval state contains \((l,u)\); an implementation may also cache denominator and height bounds.

The invariant is
\begin{equation}\label{eq:cubic-invariant}
 l\le u,
 \qquad F(l)\le x\le F(u),
 \qquad u-l=2M/2^k,
 \qquad -M\le l\le u\le M.
\end{equation}
Each clause is quantifier-free after rational comparisons are expanded to integer cross products.

\begin{lemma}[Bisection invariant]\label{lem:cubic-bisect}
The initial interval satisfies \eqref{eq:cubic-invariant}.  If a state at stage \(k\) satisfies it, the bisection transition produces a state at stage \(k+1\) satisfying it.
\end{lemma}

\begin{proof}
The base bracket was proved above.  For the step, the midpoint lies between the endpoints.  If \(F(m)\le x\), then \([m,u]\) has the required sign bracket; otherwise decidability gives \(x<F(m)\), so \([l,m]\) has the bracket.  Both new intervals lie inside the old one and have half its width.  The outer bound by \([-M,M]\) is preserved.
\end{proof}

This is exactly one application of quantifier-free induction.

\subsection{Primitive-recursive precision fuel}

Write the rational input as \(a/b\) with \(b>0\).  Since
\[
 M=\frac{|a|+b}{b}\le |a|+b,
\]
the initial width satisfies \(2M\le 2(|a|+b)\).  Let \(H=|a|+b\).  A deliberately coarse fuel is
\[
 N(q,a,b)=q+2H+3.
\]
Using \(2^n\ge n+1\), we have \(2^{2H+3}\ge 2H\) for all \(H\ge1\).  Therefore
\[
 \frac{2M}{2^{N(q,a,b)}}\le 2^{-q}.
\]
The function \(N\) is primitive recursive and monotone in \(q\).  A tighter implementation may use bit length, but tightness is irrelevant to the logical benchmark.

\begin{theorem}[Certified enclosure]\label{thm:cubic-enclosure}
There are primitive-recursive rational terms \(l(x,q),u(x,q)\) such that \iPRA\ proves
\begin{align*}
 &l(x,q)\le u(x,q),\\
 &F(l(x,q))\le x\le F(u(x,q)),\\
 &u(x,q)-l(x,q)\le2^{-q},\\
 &[l(x,q+1),u(x,q+1)]\subseteq[l(x,q),u(x,q)].
\end{align*}
\end{theorem}

\begin{proof}
Run the bisection trace for \(N(q,a,b)\) steps and apply \cref{lem:cubic-bisect}.  The width estimate follows from the fuel inequality.  Monotonicity of \(N\) and the deterministic recurrence show that the longer trace extends the shorter one, hence the intervals are nested.
\end{proof}

\subsection{Regular midpoint name}

Define
\[
 c(q)=\frac{l(x,q+2)+u(x,q+2)}2.
\]
Nestedness implies that for \(p\ge q\), both \(c(p)\) and \(c(q)\) lie in \(I_{q+2}\).  Therefore
\[
 |c(p)-c(q)|\le \operatorname{width}(I_{q+2})\le2^{-(q+2)}.
\]
By symmetry,
\[
 |c(p)-c(q)|\le2^{-(\min(p,q)+2)}
 \le2^{-p}+2^{-q}.
\]
Thus \(c\) is a regular primitive-recursive rational name.

\subsection{The equation in represented form}

To show that the represented real named by \(c\) satisfies \eqref{eq:cubic-equation}, bound the variation of \(F\) on \([-M,M]\).  For \(r,s\) in this interval,
\begin{align*}
 |F(s)-F(r)|
 &=|s-r|\,|r^2+rs+s^2+1|\\
 &\le |s-r|(3M^2+1).
\end{align*}
Let \(K\) be any primitive-recursive natural bound with
\[
 3M^2+1\le2^K.
\]
A coarse bound obtained from \(|a|+b\) suffices.  Compute the bracket at internal precision \(q+K+3\).  Since \(x\) lies between \(F(l)\) and \(F(u)\), and the midpoint lies between \(l\) and \(u\), monotonicity and the Lipschitz bound give
\[
 |F(c(q+K+3))-x|\le (3M^2+1)(u-l)\le2^{-q-3}.
\]
After aligning the regular-name convention, this proves equality of the represented reals \(F(c)=x\).

The internal precision shift is a simple demand transformer.  In this example there is no demand cycle: the input height determines \(K\), and the requested equation precision determines the bisection depth.

\subsection{Uniqueness}

Suppose regular names \(y,z\) both satisfy \(F(y)=x=F(z)\).  If they were apart by a positive rational gap \(\varepsilon\), sufficiently fine approximants would preserve an order gap, say \(y_q+\varepsilon/2<z_q\).  Strict monotonicity would then give a positive lower gap between \(F(y_q)\) and \(F(z_q)\), while equality to \(x\) makes both arbitrarily close to the same rational approximant.  This is a contradiction.  Quantitatively, the factor in \cref{lem:cubic-monotone} is at least \(1\), so
\[
 |F(s)-F(r)|\ge |s-r|
\]
for rationals \(r,s\).  Therefore equality of function values directly implies equality of represented inputs with no inverse modulus loss.

\begin{theorem}[Constructive cubic inverse]\label{thm:cubic-inverse}
For every rational \(x\), the primitive-recursive midpoint name above is the unique represented real \(y\) satisfying \(y^3+y=x\).  At precision \(q\), its certified interval is produced by a closed finite \CIFER\ graph with no residual boundary.
\end{theorem}

\begin{proof}
Existence is the enclosure, regularity, and equation argument.  Uniqueness follows from the quantitative strict monotonicity inequality.
\end{proof}

\subsection{Proof graph}

A concrete graph can use the following nodes:

\begin{longtable}{r p{0.22\textwidth} p{0.57\textwidth}}
\toprule
ID & Node & Local content \\
\midrule
\endhead
0 & input validity & denominator positive and rational code well formed \\
1 & size bound & compute \(M,H,K,N\) and prove coarse inequalities \\
2 & initial bracket & prove \(F(-M)\le x\le F(M)\) \\
3 & transition algebra & midpoint and branch equations \\
4 & invariant step & preserve bracket, nesting, and width \\
5 & trace theorem & apply quantifier-free induction for fuel \(N\) \\
6 & precision theorem & derive width \(\le2^{-q}\) \\
7 & regularity & compare nested midpoint intervals \\
8 & equation & apply the variation bound at shifted precision \\
9 & uniqueness & use inverse Lipschitz monotonicity \\
10 & final & package the represented root theorem \\
\bottomrule
\end{longtable}

Every node is local arithmetic or one bounded trace.  There is no boundary node and no infinite choice.

\subsection{Lean completion path for the benchmark}

The recommended order is:

\begin{enumerate}
\item reuse or define signed-integer and positive-denominator rational codes;
\item prove cross-multiplied equality and order reflection;
\item define \code{cubic}, midpoint, absolute-value bound, and one bisection step as \code{PR} syntax;
\item define the invariant formula and prove \code{Fm.QF};
\item prove the polynomial factorization and positivity lemmas;
\item construct base and step derivations and invoke \code{CIFER.loop};
\item prove the coarse exponential fuel inequality;
\item package interval output and regular midpoint name as fixed \code{PR} symbols;
\item prove the equation and uniqueness at the formula level;
\item expose a theorem whose type is literally \code{[] |- [PRAAx] cubicInverseFormula} and export it to \code{NPrf}.
\end{enumerate}

The supplied Python example uses exact \code{Fraction} arithmetic and checks the bracket and width invariant for several rational inputs.  These tests validate the reference transition, not the universal theorem; the object proof remains the decisive milestone.

\section{Lean formalization architecture}

\subsection{Design goal}

The formalization should make it difficult to confuse three kinds of theorem:

\begin{enumerate}
\item a host theorem about an executable function;
\item a theorem that a finite certificate verifier is correct;
\item an embedded derivation in \iPRA.
\end{enumerate}

All three are useful, but the public endpoint for an iPRA application must have type \code{Prf PRAAx}.  The architecture below organizes modules so that the transition between layers is explicit.

\subsection{Proposed module graph}

\begin{description}[style=nextline]
\item[\code{CIFER.ObjectCompiler}] Common-context proof schedules, residual discharge, bounded-loop composition, stable-observer composition, and native export.
\item[\code{CIFER.ResidualGraph}] Host graph syntax, topological validation, reachability, and residual extraction.
\item[\code{CIFER.Boolean}] Boolean terms, truth tables, fixed-formula reflection, and bounded folds.
\item[\code{CIFER.BoundedSearch}] Least witness, complete enumeration, finite extrema, and list membership reflection.
\item[\code{CIFER.Coding}] Fixed tuple, bounded sequence, finite map, and certified-setoid support.
\item[\code{CIFER.DemandClosure}] Finite vectors, monotone inflationary iteration, rank stabilization, and object trace.
\item[\code{CIFER.Relational}] Layered finite clouds, viability, least paths, stable observers, diameter, and gap collapse.
\item[\code{CIFER.Represented}] Regular rational names, intervals, extensionality, and finite-use operators.
\item[\code{CIFER.Applications.*}] Domain-specific interface instantiations.
\end{description}

The delivered artifact includes the first, second, demand, relational, and application-interface outlines.  Boolean and coding modules are listed as completion milestones because they should preferably reuse existing repository arithmetic rather than duplicate it.

\subsection{ResidualBuild}

The key structure is:

\begin{lstlisting}
structure ResidualBuild (Gamma boundary done : List (Fm PTm n)) where
  derive : forall A, A in done ->
    boundary ++ Gamma |- [PRAAx] A
\end{lstlisting}

The empty build is vacuous.  \code{step} accepts a local derivation from all prior nodes plus the permanent context; \code{stepFrom} restricts the local proof to an explicit dependency sublist.  Both use \code{Prf.trans} to replace dependency hypotheses by stored derivations.  \code{close} replaces every boundary hypothesis by a supplied proof from \(\Gamma\).

This structure is deliberately small.  It depends only on the checked deduction API, not on graph algorithms or arithmetic semantics.  It can therefore be elaborated and audited before the rest of the framework.

\subsection{Exact scoped syntax}

The upstream library indexes terms and formulas by the number of free variables.  Every module must preserve this discipline.  Recommended conventions are:

\begin{itemize}
\item use existing \code{Syn.subst}, \code{Fm.subst}, \code{Fm.shift}, \code{inst}, and \code{repl};
\item prove formula-constructor substitution lemmas once;
\item avoid unscoped integer indices in generated syntax;
\item package application substitutions as \code{Fin k -> PTm n};
\item use upstream shift/substitution identities rather than local simplifier folklore.
\end{itemize}

Generated code should keep an explicit mapping from human variable names to \code{Fin} indices in metadata.  This is not part of the proof but prevents source-audit mistakes.

\subsection{Formula Booleanization strategy}

There are two implementation routes.

\paragraph{Typed fixed syntax.}
Define a host inductive \code{BForm n} containing equations, conjunction, disjunction, implication, and strictly bounded quantifiers.  A compiler returns:
\begin{itemize}
\item a \code{PR} Boolean checker;
\item a formula in \code{Fm PTm n};
\item embedded reflection derivations.
\end{itemize}
This route provides one reusable compiler and is recommended.

\paragraph{Application-specific generation.}
Generate Boolean terms and proofs directly from application formulas.  This is simpler for a first benchmark but duplicates logical cases.

The typed route still does not create an object-level universal evaluator.  \code{BForm} recursion runs in Lean and emits fixed \code{PR} syntax.

\subsection{Representing proof-producing compilation}

A fully typed node can be represented by dependent records, but dependent arrays indexed by dynamically computed graph positions become cumbersome.  A practical two-stage design is:

\begin{enumerate}
\item validate and normalize an untyped metadata graph;
\item generate a Lean declaration for each node with its exact formula and local proof;
\item generate a topological fold expression using \code{ResidualBuild.stepFrom}.
\end{enumerate}

The generated Lean source is the proof-carrying artifact.  This avoids implementing a large dependent graph interpreter inside Lean while retaining kernel checking of every local proof and composition step.

For hand-written graphs, a typed heterogeneous list can be introduced later.  It is not required for soundness.

\subsection{Demand closure formalization}

At the host level, \code{Demand n := Fin n -> Nat}.  Prove:
\begin{align*}
 a&\le a\vee F(a),\\
 a\vee F(a)\ne a&\to \rho(a)<\rho(a\vee F(a)),\\
 a\le C&\to\rho(a)\le\rho(C).
\end{align*}
The stabilization theorem follows by natural induction.

For the object version, fixed dimension is compiled to tuple-coded natural terms.  The cap and seed are terms.  A primitive-recursive transition returns the componentwise maximum.  The invariant includes the current rank and fixed flag.  The host theorem about finite vectors guides the object proof but does not replace it.

\subsection{Relational viability formalization}

Encode a layered system by:
\begin{itemize}
\item the number of levels;
\item a length for each level;
\item an entry function \(L(k,i)\);
\item a Boolean edge checker \(E(k,i,j)\).
\end{itemize}
The backwards table is a bounded Boolean array.  Reverse time \(r=0,\ldots,h\) corresponds to level \(h-r\), making the recurrence a standard primitive recursion.

The invariant states that the table after \(r\) steps agrees with path existence to the last \(r\) levels.  Path existence is itself represented by a bounded Boolean fold or by a stored successor table.  To construct the least path, maintain the current index and use least bounded search in the next viable row.

A stable observer theorem should take object formulas proving nonemptiness and pairwise equality, then derive the observation of the least member.  Approximate diameter uses finite maxima of rational distances.

\subsection{Boundary interfaces as structures, not axioms}

The application file defines structures such as:

\begin{lstlisting}
structure HeightListInterface (Gamma : List (Fm PTm n)) where
  heightBound : Fm PTm n
  completeList : Fm PTm n
  enumerate : heightBound :: Gamma |- [PRAAx] completeList
\end{lstlisting}

This is safe because a value of the structure must contain the downstream derivation.  It does not assert that \code{heightBound} is provable.  A separate producer structure may contain a term and proof of the boundary.  Closing the theorem requires both values.

Unproved targets should be represented as definitions of formulas or propositions, comments, or data fields.  They should not be declared as \code{axiom} and should not receive a theorem with a placeholder proof.

\subsection{Application endpoint pattern}

A closed application should expose declarations resembling:

\begin{lstlisting}
def graph : ... := ...

theorem residual_result :
    boundary ++ Gamma |- [PRAAx] goal := by
  -- generated topological fold

theorem producer_1 : Gamma |- [PRAAx] boundary[0] := ...
...

theorem closed_result : Gamma |- [PRAAx] goal :=
  residualBuild.close ...

def native_result : NPrf ... :=
  pra_iff_native.mp closed_result
\end{lstlisting}

The separate names make status transparent.  During development, \code{residual\_result} may compile before the producer theorems.  The public README can then report exactly which layer is complete.

\subsection{Automation boundaries}

Tactics may assist with:
\begin{itemize}
\item formula substitution simplification;
\item propositional natural-deduction plumbing;
\item ring-normalization certificate generation;
\item bounded arithmetic normalization;
\item construction of repetitive proof-graph folds.
\end{itemize}

Every tactic should produce ordinary proof terms checked by Lean.  The framework should not add an unverified oracle.  Expensive external algebra systems may produce finite certificates -- polynomial identities, factorizations, interval inequalities -- which small verified checkers consume.

\subsection{Testing strategy}

Use four independent layers of tests.

\begin{enumerate}
\item \textbf{Reference tests:} exact Python computations for graph algorithms and numerical examples.
\item \textbf{Native Lean tests:} finite algorithms and semantic invariants over Lean data types.
\item \textbf{Embedded unit proofs:} small \code{Prf PRAAx} examples for each generic combinator.
\item \textbf{End-to-end endpoint:} at least one nontrivial closed application with a printed axiom report.
\end{enumerate}

Passing the first two does not imply the fourth.  They are debugging and regression layers.

\subsection{Build and audit workflow}

The pinned repository toolchain is Lean 4.32.0.  A formalization run should:

\begin{enumerate}
\item place \code{CIFER} on the same import path as \code{ConciseIPRA};
\item build the composition core first;
\item record compiler output and hashes;
\item run source searches for admissions and new axioms;
\item use \code{\#print axioms} on public endpoints;
\item save the exact theorem type;
\item regenerate the residual and graph manifests;
\item compare finite reference output with object-level numeral instances where feasible.
\end{enumerate}

A generated package should keep build logs out of the mathematical source tree or clearly mark them as environment-specific.

\subsection{Suggested implementation order}

The following order minimizes risk.

\begin{enumerate}
\item Elaborate \code{ObjectCompiler.lean} against current upstream APIs.
\item Close a propositional three-node residual example.
\item Add one natural-number loop and prove a closed bounded-search instance.
\item Implement Boolean reflection and generic bounded search.
\item Implement rational coding and the cubic benchmark.
\item Prove residual graph soundness and generate a proof schedule automatically.
\item Add finite viability and stable observers.
\item Add capped demand closure only after a real application needs a cycle.
\item Migrate height/list and gap-conditional application interfaces.
\end{enumerate}

This order ensures that the actual embedded endpoint is tested before the framework grows broad.

\subsection{Expected proof bottlenecks}

The most likely engineering bottlenecks are not the mathematical metatheorems.  They are:

\begin{itemize}
\item scoped substitution under bounded quantifier encodings;
\item reusable arithmetic reflection for signed and rational codes;
\item keeping invariants syntactically quantifier-free;
\item proof-term size for generated finite folds;
\item normalization of contexts and list membership proofs;
\item connecting existing native arithmetic lemmas to newly built embedded derivations without assuming a reflection principle that has not been proved.
\end{itemize}

Each bottleneck has a local remedy: stronger syntax utilities, proof-generating reflection lemmas, Boolean accumulators, shared definitions, context combinators, and explicit object proofs.

\section{Limits, research program, and conclusions}

\subsection{What is proved at the human level}

The generic arguments in this report are finitary and constructive.  They establish:

\begin{itemize}
\item residual compilation and boundary discharge for fixed finite proof DAGs;
\item Boolean reflection for fixed quantifier-free and strictly bounded formulas;
\item least bounded search and complete bounded enumeration;
\item capped monotone demand closure with an explicit rank fuel;
\item finite-trace reflection by quantifier-free induction;
\item backward viability and least paths at fixed finite horizons;
\item exact and approximate stable-observer elimination;
\item finite-horizon Cauchy coherence under explicit shrinking and motion moduli;
\item height/list conversion for decidable coded predicates;
\item a closed rational bisection construction for the inverse of \(y^3+y\).
\end{itemize}

These are mathematical theorem schemas for every fixed collection of primitive-recursive symbols satisfying the stated local premises.  The proofs do not rely on excluded middle for arbitrary formulas or on countable choice.

\subsection{What remains to be formalized}

The Lean package is an integration-oriented outline.  The object composition core is written against the inspected API, but this artifact environment did not contain Lean or Lake, so it was not elaborated.  The bounded-formula compiler, rational arithmetic reflection, object demand trace, relational viability proof, and cubic endpoint remain implementation milestones.

Accordingly, the report makes no new kernel-checking claim.  The first decisive formal result should be the cubic benchmark ending in a declaration of type \code{Prf PRAAx}.  Until that endpoint builds, the framework should be described as a detailed formalization plan with source-level combinators, not as a completed embedded library.

\subsection{What a general framework cannot manufacture}

Residualization is not a mechanism for proving arbitrary missing theorems.  It exposes them.  The following remain theorem-specific:

\begin{itemize}
\item effective height bounds for general curve families where none is presently supplied;
\item rational representatives for general Mordell--Weil quotient data;
\item explicit analytic constants for strong uniform prime-distribution estimates;
\item positive separation gaps in collision-prone exact-selection problems;
\item triangulation, cover, or compactness moduli not derivable from the given representation;
\item convergence rates for asymptotic theorems stated without quantitative data.
\end{itemize}

The framework may still improve these research problems by giving each one a precise input/output type and by proving every elementary consequence once the input is available.

\subsection{No generic logical minimality claim}

The residual list is canonical under graph reachability, duplicate removal, and declared internal derivations.  It is not guaranteed to be a smallest logically sufficient subset of formulas.  Such minimization would require solving arbitrary theorem-proving problems and would be unstable under reformulation.

This limitation is acceptable.  Dependency-normal residuals are reproducible and explainable.  Domain experts may add explicit implication nodes to eliminate dominated assumptions.

\subsection{No internal universal compiler}

The host software may accept many graph descriptions, but each object proof is specialized to one finite graph.  The report does not assert a primitive-recursive evaluator for arbitrary \code{PR} syntax, an internal proof compiler for arbitrary encoded derivations, or a reflection theorem saying that successful host computation implies object derivability.

This boundary is required by the computational strength of PRA and by the design of a small trusted core.

\subsection{Potential improvements}

Several refinements are promising after the first kernel-checked benchmark.

\paragraph{Proof-preserving graph optimization.}
Fuse adjacent deterministic nodes, share repeated subproofs, and remove unused formulas while preserving a generated equivalence proof between original and optimized schedules.

\paragraph{Residual equivalence libraries.}
Collect embedded implications between common boundary forms: modulus strengths, list and height interfaces, alternative norm bounds, and gap certificates.  Residual normalization can then use proved dominance rather than heuristics.

\paragraph{Certified code generation.}
Define a small host IR whose compiler emits both \code{PR} syntax and local proof templates.  Prove the compiler correct once at the host theorem level, while continuing to kernel-check every emitted object derivation.

\paragraph{Domain-specific invariant combinators.}
Provide libraries for scan invariants, finite extrema, frozen descent, interval refinement, dynamic programming, and table replay.  These should reduce application proofs to local arithmetic obligations.

\paragraph{Proof-cost accounting.}
Record graph size, number of cuts, invariant size, and generated proof-term size.  This would make the ``simplest proof checker implementation'' objective measurable alongside mathematical generality.

\paragraph{Residual-driven research planning.}
Aggregate residual reports across packets.  A producer that discharges many high-impact residuals should receive priority.  This turns the repository into a dependency map of open quantitative lemmas.

\subsection{Recommended repository integration}

The new framework should be integrated as a small library rather than another isolated packet.

\begin{enumerate}
\item Put the generic object compiler under \code{ConciseIPRA} or a sibling checked directory after it elaborates.
\item Keep research reports and generated examples under \code{docs/proofs/foundations}.
\item Give every application packet a machine-readable boundary manifest.
\item Add a CI check that distinguishes native Lean theorems from \code{Prf PRAAx} endpoints.
\item Maintain one catalog table with columns for human proof, reference tests, native Lean, embedded Lean, residuals, and upstream commit.
\end{enumerate}

This would reduce repeated status prose and make it easier to see which generic infrastructure is truly shared.

\subsection{Further research questions}

\begin{enumerate}
\item Can a small, checked collection of proof-producing invariant combinators cover most current packets without a large tactic framework?
\item Which existing arithmetic modules already contain enough embedded rational and bounded-list infrastructure for the cubic benchmark, and where is native-only duplication occurring?
\item Can generated proof graphs be normalized to a compact sequent calculus representation and then translated to the current natural-deduction \code{Prf} with a checked compiler?
\item What is the best finite relational interface for polynomial roots that remains stable at multiplicity changes and supports common algebraic consumers?
\item Which Diophantine packets share the same height-bound residual after canonicalization of curve and point codes?
\item Can explicit ``count plus isolation oracle'' assumptions be factored into a weaker reusable boundary than a full height bound?
\item Which analytic number theory estimates have effective constants that are elementary enough to close within PRA, and which require stronger constructive systems?
\item Can finite-horizon coherence be generalized from metric observations to certified uniform spaces or apartness spaces while remaining first-order arithmetic after coding?
\item How much proof-term compression is achieved by one fused trace versus a modular proof DAG on representative applications?
\item Can the residual graph drive automatic generation of theorem statements, READMEs, and status ledgers without weakening human review?
\end{enumerate}

\subsection{Conclusion}

The main lesson is methodological.  Mathematics in \iPRA\ should not begin with an ideal object and ask how to force all of it into one numerical sort.  It should begin with the finite observation required by one theorem.  Exact computations become primitive-recursive terms; bounded processes become traces with quantifier-free invariants; nonunique data become finite relations; stable consumers eliminate unnecessary choice; and genuinely missing quantitative producers remain visible as residual formulas.

Under this discipline, a general framework can be both ambitious and honest.  It can compile large classes of finite mathematics to actual object-language derivations, reuse those derivations across domains, and still state precisely where current mathematics has not supplied an effective bound or selector.  The resulting endpoint is not merely an executable program or a semantic Lean theorem.  It is a sequent of \iPRA, closed exactly when its producer obligations have been proved.

\appendix
\section{Theorem and implementation ledger}

The status terms used below are:

\begin{description}[style=nextline]
\item[Human-complete.] A full mathematical argument is given in this report.
\item[Finite-tested.] The exact Python reference exercises nontrivial finite instances.
\item[Lean source.] Definitions or proof-composition source is supplied, but no build is claimed.
\item[Formalization target.] A precise endpoint and dependency path are specified, but the embedded proof is not yet constructed.
\item[Residual/open.] A theorem-specific producer is deliberately left as a boundary.
\end{description}

\small
\begin{longtable}{p{0.07\textwidth}p{0.34\textwidth}p{0.16\textwidth}p{0.33\textwidth}}
\toprule
ID & Claim & Status & Formal route \\
\midrule
\endfirsthead
\toprule
ID & Claim & Status & Formal route \\
\midrule
\endhead
R1 & Reachability marks exactly the dependency closure of a target in a finite topological graph. & Human-complete; finite-tested & \code{ResidualGraph.reachable\_correct}. \\
R2 & Reachable primitive boundary leaves form the canonical dependency-normal residual. & Human-complete; finite-tested & \code{residual\_characterization}. \\
R3 & Every graph node is derivable from its residual boundary plus the permanent context. & Human-complete; Lean source & Fold with \code{ResidualBuild.stepFrom}. \\
R4 & Supplying embedded proofs of all residual formulas closes the target. & Human-complete; Lean source & \code{ResidualBuild.close}; simultaneous cut. \\
R5 & Improving a producer substitutes its residual into all downstream residuals without changing downstream local proofs. & Human-complete & Recursive residual equation and context weakening. \\
B1 & Every fixed quantifier-free formula has a primitive-recursive Booleanizer with two-way reflection. & Human-complete; formalization target & Host recursion on fixed syntax; equality reflection and truth tables. \\
B2 & Every fixed strictly bounded formula has a reflected Boolean fold. & Human-complete; formalization target & \code{PRA.qfInd} on conjunction/disjunction folds. \\
S1 & Least bounded search proves soundness, minimality, and the no-witness branch. & Human-complete; finite-tested pattern & Search-state invariant; \code{CIFER.loop}. \\
S2 & Bounded filtering gives a sound, complete, increasing witness list. & Human-complete; formalization target & Sequence state and bounded completeness invariant. \\
S3 & Finite maxima/minima and least attaining indices are primitive recursive and reflected. & Human-complete; finite-tested use & Scan invariant. \\
D1 & Capped inflationary demand iteration stabilizes within the rank bound. & Human-complete; finite-tested; Lean source definitions & Complete proofs in \code{DemandClosure.lean}; object trace after tuple coding. \\
D2 & Under monotonicity, the stabilized demand is the least prefixed point above the seed. & Human-complete & Induction comparing every iterate with an arbitrary prefixed point. \\
T1 & A fixed bounded trace with QF invariant yields its invariant at arbitrary PR fuel. & Human-complete; Lean source & \code{invariantAt} from \code{PRA.qfInd}. \\
T2 & A terminal extractor may conclude an arbitrary quantified formula. & Human-complete; Lean source & \code{CIFER.loop}; implication elimination. \\
T3 & Freezing after termination turns partial early exit into a fixed-fuel primitive recursion. & Human-complete; formalization target & Boolean done flag in invariant. \\
P1 & Topological proof schedules compose local embedded derivations by simultaneous cut. & Human-complete; Lean source & \code{ResidualBuild.step} and \code{stepFrom}. \\
V1 & Backward viability is equivalent to existence of a finite compatible path. & Human-complete; finite-tested; Lean source definitions & Reverse-time Boolean table; induction on horizon. \\
V2 & A nonempty viable system has a lexicographically least finite path. & Human-complete; finite-tested & Repeated least bounded search. \\
O1 & Exact stable observations eliminate finite witness choice. & Human-complete; finite-tested; Lean composition source & Least viable member plus pairwise stability. \\
O2 & Diameter-bounded observations determine an approximate branch-independent output. & Human-complete; finite-tested & Finite maximum of rational pairwise distances. \\
O3 & A positive separation gap collapses an approximate cloud to one exact label. & Human-complete; formalization target & Triangle inequality plus finite label search. \\
C1 & Changing finite horizons yield a Cauchy name under diameter and motion-tail moduli. & Human-complete; formalization target & Common-horizon comparison; no infinite branch choice. \\
H1 & A positive height bound yields a complete finite list for a decidable coded predicate. & Human-complete & Bounded code enumeration and filter reflection. \\
H2 & A complete finite list yields a height bound by finite maximum. & Human-complete & List completeness plus maximum theorem. \\
H3 & A cardinality bound alone is not a search cutoff or complete list producer. & Logical distinction explained & Separate interface formulas; no conversion node. \\
A1 & Finite auxiliary exhaustion avoids a global selector. & Human-complete & Bounded union of branch outputs plus target filter. \\
E1 & Exact rational bisection preserves the bracket for \(y^3+y=x\). & Human-complete; finite-tested & Rational order reflection and QF trace invariant. \\
E2 & Coarse primitive-recursive fuel gives interval width at most \(2^{-q}\). & Human-complete; finite-tested & Exponential lower bound and trace width equation. \\
E3 & Midpoints form a regular name satisfying the cubic equation and uniqueness. & Human-complete; formalization target & Nested intervals, Lipschitz bound, strict monotonicity. \\
X1 & General effective Faltings/Siegel height producer. & Residual/open & Supply an embedded \code{HeightBound} proof; downstream H1 closes. \\
X2 & General rational Mordell--Weil covering representatives. & Residual/open & Supply \code{RationalCover}; bounded descent then closes. \\
X3 & Full effective constants for strong uniform prime-distribution estimates. & Residual/open & Supply \code{EffectiveEstimate}; finite arithmetic consumer closes. \\
X4 & Exact selector through arbitrary collisions without a gap. & Deliberately not claimed & Use a finite relation/stable observable, or supply a positive gap. \\
\bottomrule
\end{longtable}
\normalsize

\section{Application migration checklist}

For each proposed theorem, record the following items before writing a large proof script.

\begin{enumerate}
\item \textbf{Target formula.} State the exact \code{Fm PTm n} endpoint, including free parameters and guards.
\item \textbf{Observation.} Identify the finite output actually consumed: Boolean, rational interval, list, finite relation, equality certificate, or represented name.
\item \textbf{Codes.} Fix canonical or setoid codes and all validity predicates.
\item \textbf{Producer list.} Name every algorithm or theorem that must output finite evidence.
\item \textbf{Boundary list.} Mark producers not yet supplied by actual embedded derivations.
\item \textbf{Demand map.} Give output-to-input precision and resource transformers, including caps for cycles.
\item \textbf{Trace list.} Give each bounded loop a state, transition, fuel, and explicit QF invariant.
\item \textbf{Local proof list.} Give each proof node an actual dependency sequent.
\item \textbf{Relational choice audit.} Explain whether a selector is finite, stable, gap-guarded, or unnecessary.
\item \textbf{Endpoint audit.} Confirm whether the public declaration is native Lean, \code{Prf PRAAx}, or \code{NPrf}.
\item \textbf{Residual report.} Export reachable boundary formulas and remove irrelevant ones.
\item \textbf{Validation report.} Separate finite tests, PDF/human proof review, Lean build, and axiom audit.
\end{enumerate}

\section{Source and provenance notes}

The repository audit was performed against the public \code{main} branch of
\url{https://github.com/VladimirReshetnikov/ProveIt}, with current head
\code{8649990597d73657bde480b6c389de7f85bc0cb0} observed on 15 November 2026 in America/Los Angeles.  The principal inspected paths were:

\begin{itemize}
\item \code{Logic/PrimitiveRecursiveArithmetic/README.md};
\item \code{Logic/PrimitiveRecursiveArithmetic/ConciseIPRA/Lean/ConciseIPRA/}\newline
      \code{Syntax.lean}, \code{Deduction.lean}, \code{PR.lean}, and \code{NativeBiInterpretation.lean};
\item \code{Logic/PrimitiveRecursiveArithmetic/docs/proofs/README.md};
\item foundation packets for proof-carrying observation graphs, OCiPRA, PFIG, unified finitary methods, and kernel finite-use compilation;
\item difficulty reports for Faltings, Siegel, Mordell--Weil, Euler--Rodrigues, and prime distribution in \iPRA.
\end{itemize}

The report treats repository statements about packet compilation status as status records, not as independent kernel verification.  No upstream file was modified by this artifact.

\begin{thebibliography}{99}

\bibitem{BishopBridges1985}
E. Bishop and D. Bridges,
\emph{Constructive Analysis},
Grundlehren der mathematischen Wissenschaften 279,
Springer, 1985.
DOI: \url{https://doi.org/10.1007/978-3-642-61667-9}.

\bibitem{BridgesRichman1987}
D. Bridges and F. Richman,
\emph{Varieties of Constructive Mathematics},
London Mathematical Society Lecture Note Series 97,
Cambridge University Press, 1987.

\bibitem{deMouraEtAl2015}
L. de Moura, S. Kong, J. Avigad, F. van Doorn, and J. von Raumer,
``The Lean Theorem Prover (System Description),''
in \emph{Automated Deduction -- CADE-25}, LNCS 9195,
pp. 378--388, Springer, 2015.
DOI: \url{https://doi.org/10.1007/978-3-319-21401-6_26}.

\bibitem{Godel1958}
K. G\"odel,
``\"Uber eine bisher noch nicht ben\"utzte Erweiterung des finiten Standpunktes,''
\emph{Dialectica} 12 (1958), 280--287.
\newline DOI: \nolinkurl{10.1111/j.1746-8361.1958.tb01464.x}.

\bibitem{Kleene1952}
S. C. Kleene,
\emph{Introduction to Metamathematics},
North-Holland, 1952.

\bibitem{Kohlenbach2008}
U. Kohlenbach,
\emph{Applied Proof Theory: Proof Interpretations and their Use in Mathematics},
Springer Monographs in Mathematics, Springer, 2008.
DOI: \url{https://doi.org/10.1007/978-3-540-77533-1}.

\bibitem{TroelstraVanDalen1988}
A. S. Troelstra and D. van Dalen,
\emph{Constructivism in Mathematics: An Introduction}, Volumes I and II,
North-Holland, 1988.

\bibitem{ProveIt}
V. Reshetnikov,
\emph{ProveIt repository},
\url{https://github.com/VladimirReshetnikov/ProveIt},
audited at the commit recorded above.

\end{thebibliography}


\end{document}
